\documentclass[UTF8,11pt]{ctexart}
\usepackage{amsmath,amssymb,amsthm,mathtools}
\usepackage{booktabs,longtable,array,geometry,enumitem}
\usepackage{xcolor}
\definecolor{linkblue}{rgb}{0.10,0.20,0.45}
\usepackage[colorlinks=true,linkcolor=linkblue,citecolor=linkblue,urlcolor=linkblue]{hyperref}
\geometry{a4paper,margin=2.3cm}
\setlength{\parskip}{3pt}
\setlength{\emergencystretch}{3em}
\sloppy
\newtheorem{theorem}{定理}[section]
\newtheorem{lemma}[theorem]{引理}
\newtheorem{proposition}[theorem]{命题}
\newtheorem{corollary}[theorem]{推论}
\theoremstyle{definition}
\newtheorem{remark}[theorem]{注}
\newcommand{\Lam}{\Lambda}
\newcommand{\tr}{\operatorname{tr}}
\newcommand{\dd}{\,\mathrm{d}}
\newcommand{\R}{\mathbb{R}}
\newcommand{\E}{\mathbb{E}}
\newcommand{\one}{\mathbf{1}}
\newcommand{\esssup}{\operatorname*{ess\,sup}}
\newcommand{\essinf}{\operatorname*{ess\,inf}}
\newcommand{\Var}{\operatorname{Var}}
\newcommand{\diag}{\operatorname{diag}}
\newcommand{\Sym}{\operatorname{Sym}}
% \peq: \u5f15\u7528\u53d7\u5ba1\u7a3f\u4ef6\u7684\u516c\u5f0f\uff08\u7528\u5176 \label \u540d\uff0c\u4e0d\u4ea7\u751f\u672c\u6587\u672a\u5b9a\u4e49\u5f15\u7528\uff09
\newcommand{\peq}[1]{\textup{\u5f0f}\,\texttt{#1}}
\newcommand{\ok}{\textcolor[rgb]{0.05,0.45,0.15}{\bfseries [已证]}}
\newcommand{\chk}{\textcolor[rgb]{0.05,0.25,0.55}{\bfseries [已核]}}
\newcommand{\gap}{\textcolor[rgb]{0.70,0.15,0.10}{\bfseries [缺口]}}
\newcommand{\num}{\textcolor[rgb]{0.55,0.35,0.00}{\bfseries [数值]}}
\title{\bfseries 稿件审核报告\\[4pt]
\large \emph{A joint moment bound for the three largest graph eigenvalues}\\[2pt]
\normalsize 主定理证明链的逐节复核、符号与数值验证、问题清单及对学位论文主线的影响}
\author{}
\date{2026 年 11 月 2 日}
\begin{document}
\maketitle
\vspace{-1.5em}

\begin{abstract}
本报告对稿件 \emph{A joint moment bound for the three largest graph eigenvalues}
（\texttt{top3\_joint\_moment\_bound.pdf}，19 页；权威源文件
\texttt{top3\_joint\_moment\_bound.tex}，1346 行）进行独立审核。稿件主张：对任意
$n\ge3$ 的有限简单图 $G$，
\[
\begin{split}
 \Lam_3(G)&=\lambda_1(G)+\lambda_2(G)+\lambda_3(G)\ \le\ c\,n-3,\\
 c&=\frac{106106902139519}{107310156800000}+\frac{41\sqrt3}{240}
    -\frac{53}{9000000}=1.284673255723918505\ldots ,
\end{split}
\]
从而把 Sivashankar--Tang--Wakhare（arXiv:2608.02429，命题 6.1）的系数
$(3+\sqrt5)/4=1.309016994\ldots$ 降低 $0.024343738651028919\ldots$；并主张每个取到极大值的
graphon 其典范律至多有 $24$ 个支撑点。

审核采用三条互相独立的轨道：(i) 对全部 $34$ 个编号结果逐行手推重证，不把稿件的中间步骤当作前提；
(ii) 用 \texttt{sympy} 精确有理数与二次域算术机检约 $120$ 条显示恒等式、因式分解与有理数证书；
(iii) 用 \texttt{numpy/scipy} 做 $16$ 条证书恒等式的残差检验、$4000$ 个随机各向同性律的压力测试、
$k=4$ 至 $k=20$ 共 $11$ 个原子数下的全局极值搜索（每个 $60$--$70$ 组随机初值）、$n=3$ 至 $6$ 全部
$33864$ 个图的穷举，以及 $n$ 至 $2016$ 的吹起族验证。

\textbf{总体结论}：\emph{主定理的证明在数学上成立}，我没有发现任何会导致主定理失效的错误。
发现 \textbf{两处可修补的论证跳步}（\S\ref{sec:g1} 的阈值公式唯一性、\S\ref{sec:g2} 的电路论证中
"所有观测均为常值"这一情形），二者都属于表述层面而非数学错误，本报告给出完整修补（分别约 4 行与 1 句）。
其余为 $6$ 条表述性建议与 $1$ 条文献残余风险。

\textbf{必须向导师明确指出的实质性事实}：常数 $c$ \emph{不是}最优值。稿件自身构造给出下界
$c_{\mathrm w}=\rho_0=(9+\sqrt5+2\sqrt{35})/18=1.281568196872\ldots$，与 $c$ 相差
$3.105058851751\times10^{-3}$；本报告用 $n=126,\ldots,2016$ 的吹起族数值确认
$\Lam_3(G)/n\to\rho_0$，即该缺口是真实存在的、随 $n$ 线性增长的。稿件在"Further questions"
一节对此已作诚实声明。这对学位论文"求精确上界"的定位有直接影响，详见 \S\ref{sec:thesis}。
\end{abstract}

\tableofcontents
\newpage

\section{审核范围、材料与判定标准}\label{sec:scope}

\subsection{受审材料}
\begin{center}
\begin{tabular}{lll}
\toprule
文件 & 说明 & 用途\\
\midrule
\texttt{D:\textbackslash top3\_joint\_moment\_bound.pdf} & 19 页，pdfTeX & 主审对象\\
\texttt{D:\textbackslash top3\_joint\_moment\_bound.tex} & 1346 行 & \textbf{权威源}（公式一律以此为准）\\
\texttt{D:\textbackslash top3\_submission\_sources.zip} & 投稿包 & 完整性核对\\
\texttt{D:\textbackslash contribution\_and\_checks\_zh.tex} & 作者自查（中文） & 与本报告交叉比对\\
\bottomrule
\end{tabular}
\end{center}

\noindent\textbf{关于公式来源的说明。} PDF 的文本抽取会把数学式严重错排（例如
$41\sqrt3$ 被抽成 ``$\sqrt{41}\cdot3$''，$(3+\sqrt5)/4$ 被拆行），因此本报告所有公式核对
一律直接读 \texttt{.tex} 源文件。这是审核方法上的第一步，也是必要的一步。

\subsection{判定标签}
本报告对每一条结论使用四个互斥标签，避免把"我验证过"与"我相信它"混为一谈：
\begin{itemize}[leftmargin=1.6em,itemsep=1pt]
\item \ok\ 我独立重推了完整证明，每一步都成立（含精确算术）。
\item \chk\ 我用符号计算或数值计算核对为真，但没有独立写出人类可读的完整证明（例如某些
$41503$ 量级的整数系数：我验证了它们满足所述性质，但没有复现其搜索过程）。
\item \num\ 仅有数值证据。\textbf{数值证据不等同于结果}。
\item \gap\ 论证存在跳步或前提缺失。若可修补，我给出修补并说明修补后的判定。
\end{itemize}

\subsection{总体结论一览}
\begin{center}
\small
\begin{tabular}{@{}p{0.50\textwidth}p{0.42\textwidth}@{}}
\toprule
稿件主张 & 判定\\
\midrule
定理 1.1（\texttt{main}）：$\Lam_3(G)\le cn-3$ & \ok\ （依赖 G1 修补）\\
定理 1.1（\texttt{main}）：律界 $\E|X\!\cdot\!X'|+|\E X|^2\le 2c$ & \ok\\
定理 1.2（\texttt{structure-main}）：$\le24$ 支撑点 & \ok\ （依赖 G1、G2 修补）\\
定理 1.2：原 graphon 几乎处处等于二值阶梯核 & \gap\ $\to$\ok\ （G1，修补 4 行）\\
引理 2.2（\texttt{ranktwo}）：$\Sigma_2(W)\le7/6$ & \ok\\
引理 2.3（\texttt{attainment}）：$\mathcal M$ 可达，$c_{\mathrm w}\le\mathcal M\le(1+\sqrt3)/2$ & \ok\\
引理 2.4（\texttt{ellipse-structure}）：$F^TLF=\mathcal M$ a.e. & \ok\\
引理 2.5（\texttt{row-insertion}）：$r^TL^{-1}r\le\mathcal M$ & \ok\ （独立重推）\\
引理 2.6（\texttt{contact-structure}）：阈值公式 & \gap\ $\to$\ok\ （G1）\\
定理 3.1（\texttt{negative}）：二次变差半负定 & \ok\\
引理 3.2（\texttt{propagate}）：$r\le24$ & \ok\\
引理 3.3（\texttt{descent}）：$|T|\le2^{d-1}$ & \ok\\
定理 3.4（\texttt{finite}）：$|S|\le1204$ & \ok\\
引理 4.2（\texttt{binary}）：分离观测 & \gap\ $\to$\ok\ （G2，修补 1 句）\\
定理 4.3（\texttt{six}）／推论 4.4（\texttt{twentyfour}） & \ok\\
式 (5.1)--(6.4)：矩预算 \texttt{budget4}/\texttt{budget7} & \ok\ $+$\chk\\
引理 6.1（\texttt{scalar}）：$\sigma$ 的两条点态不等式 & \ok\\
引理 6.3（\texttt{radialpayment}） & \ok\\
命题 7.2（\texttt{conditional-law}）：$D\ge d_0$ & \ok\\
引理 8.2（\texttt{contact-square}）：$z_*^2\le3\sigma/2$ & \ok\ （常数 $3/2$ 非紧）\\
命题 8.5（\texttt{conditional8}）：$D\ge d_1$ & \ok\\
引理 9.1（\texttt{spectral-reserve}） & \ok\\
命题 9.3（\texttt{final-certificate}）：$\mathcal M\le c$ & \ok\\
"系数 $c$ 最优"／"极值图已确定" & 稿件\textbf{未}主张（正确）\\
\bottomrule
\end{tabular}
\end{center}

\section{稿件主张的精确复述与常数表}\label{sec:claims}

稿件的逻辑骨架可以压缩成一条链。设 $\mathcal M=\sup_W\Sigma_3(W)$（$W$ 遍历所有 graphon，
$\Sigma_3$ 为前三正特征值之和、不足补零），则
\[
 \underbrace{\mathcal M\ \ge\ c_{\mathrm w}=\rho_0}_{\text{引理 2.3：}K_1\vee C_5\text{ 加权构造}}
 \qquad\text{且}\qquad
 \underbrace{\mathcal M\ \le\ c}_{\text{命题 9.3：证书}},
\]
再由"任意各向同性律 $\to$ 二值核 $W_X$ $\to$ $\Sigma_3(W_X)\le\mathcal M$"与
"秩三 Ky Fan 公式 $\to$ 去对角"两步转移，得到有限图界 $\Lam_3(G)\le cn-3$。

\begin{center}
\small
\begin{tabular}{llll}
\toprule
量 & 精确值 & 十进制 & 来源\\
\midrule
STW 系数 & $(3+\sqrt5)/4$ & $1.309016994374947$ & \cite{STW} arXiv:2608.02429 命题 6.1\\
稿件系数 $c$ & $\frac{106106902139519}{107310156800000}+\frac{41\sqrt3}{240}-\frac{53}{9000000}$ & $1.284673255723918505$ & 定理 1.1\\
改进量 & $(3+\sqrt5)/4-c$ & $0.024343738651028919$ & 已核，与稿件一致\\
下界 $c_{\mathrm w}=\rho_0$ & $(9+\sqrt5+2\sqrt{35})/18$ & $1.281568196872168$ & 引理 2.3\\
\textbf{缺口} & $c-\rho_0$ & $\mathbf{0.003105058851751}$ & \textbf{本报告核算}\\
$d$ & $(3+\sqrt5)/2-2c$ & $0.048687477302$ & 式 (5.1)\\
$d_1$ & $d-1/3000000$ & $0.048687143969$ & 式 (5.1)\\
$d_0$ & $d_1-103/9000000$ & $0.048675699524$ & 式 (5.1)\\
$\delta_{\mathrm q}$ & $(3+\sqrt5)/2-544886/207515$ & $0.048694\ldots$ & \S6\\
分支阈值 & $1/20$ & $0.05$ & \S7--\S9\\
\bottomrule
\end{tabular}
\end{center}

\noindent\textbf{一个容易被忽略的结构性事实。} 由于 $c\in\mathbb Q(\sqrt3)$（$\sqrt5$ 在
$d=(3+\sqrt5)/2-2c$ 中被完全消去），我已用符号计算验证
$c=(3+\sqrt5)/4-d_1/2-1/6000000$，即"自对原子项"对 $c$ 的贡献恰为 $1/6000000=1.67\times10^{-7}$，
与稿件"Most of the improvement \dots\ comes from the preceding joint moment estimate"的说法一致。
换言之，\textbf{改进量的 $99.995\%$ 来自 \S5--\S8 的联合矩证书，末端的原子项几乎可以忽略}。
这一点对后续优化方向的排序很重要（见 \S\ref{sec:slack}）。

\section{审核方法与工作量}\label{sec:method}

\subsection{三轨制}
\textbf{轨道一（逐行手推）。} 对稿件 $34$ 个编号结果，我不使用其证明中的任何中间断言作为前提，
而是从定义出发重新推导。凡涉及不等式放缩处，我逐条检查放缩方向与所用常数。
特别地，对两处最容易出错的环节——引理 2.5（行插入）与引理 2.3（$\mathcal M$ 的可达性）——
我做了完全独立的重新构造，得到的表达式与稿件逐项吻合。

\textbf{轨道二（符号机检）。} 全部有理数以 \texttt{sympy.Rational} 精确表示，$\sqrt3,\sqrt5$
保留在二次域中，不使用浮点。共验证约 $120$ 条判据：因式分解恒等式、配方恒等式、系数归并关系、
有理数下界证书、以及所有"$>0$"断言。对稿件用 $\sqrt5<161/72$、$\sqrt3>433/250$ 得到的四个有理数下界，
我逐条重算，结果见 \S\ref{sec:sym}。

\textbf{轨道三（数值机检）。} 见 \S\ref{sec:num}。

\subsection{工作量}
\begin{center}
\small
\begin{tabular}{lrl}
\toprule
项目 & 数量 & 说明\\
\midrule
稿件编号结果逐行复核 & 34 & 引理/定理/命题/推论\\
稿件显示公式逐条复核 & 约 60 & 含所有编号方程\\
符号判据（精确算术） & 约 120 & 全部通过\\
证书恒等式残差检验 & 16 & 机器精度（$\le1.6\times10^{-14}$）\\
随机各向同性律压力测试 & 4000 & 违例 $0$ 例\\
全局极值搜索 & $k=4..20$ 共 11 个原子数 & 每档 $60$--$70$ 组初值\\
有限图穷举 & 33864 个（$n=3..6$） & 违例 $0$ 例\\
结构图族与随机图 & 18 个（$n\le600$） & 违例 $0$ 例\\
吹起族收敛验证 & $n=126..2016$ & $\Lam_3/n\to\rho_0$\\
\bottomrule
\end{tabular}
\end{center}

\subsection{审核局限（必须说明）}
\begin{enumerate}[leftmargin=1.8em,itemsep=2pt]
\item 对稿件中出现的大整数（如 $41503=7^3\cdot11^2$、$106106902139519$、$725029049$ 等），
我验证的是\emph{它们满足稿件所述的性质}，而不是复现其搜索/优化过程。因此我判定为 \chk\ 而非 \ok。
\item Ebrahimi 2008 年 SFU 硕士论文（定理 3.4.1）仍无法获取，方法重叠风险未能完全排除
（作者自查文件中已记录同一事项）。本报告 \S\ref{sec:novel} 给出我基于 STW 原文的间接判断：风险较低。
\item 数值验证不能替代证明。凡 \num\ 标签处，我都明确写出它不构成结果。
\item 我没有联网检索 2026 年 10 月之后的新预印本。
\end{enumerate}

\section{逐节复核}\label{sec:detail}

本节按稿件章节顺序给出复核记录。对每一条结果，我写明：它在证明链中的作用、我的复核路径、判定，
以及我认为导师会关心的评注。凡是关键的不等式或恒等式，我把自己的重推过程写出来，而不是只写"已核对"，
这样导师可以直接检查我的检查。

\subsection{引理 2.2（\texttt{ranktwo}）：$\Sigma_2(W)\le7/6$}

\textbf{作用。} 这条引理是"下界常数 $\rho_0$ 有意义"的前提：它给出
$\lambda_3=\mathcal M-\Sigma_2(W)\ge\rho_0-7/6>0$，从而每个极大 graphon 恰有三个正特征值，
后面全部的椭圆面结构与二阶变分才有对象。判定 \ok。

\textbf{我的重推。} 记 $Y=RU$（$R=|Y|$，$U$ 为单位向量，零点处任取），$C=\E R$，$N=\E[RUU^T]$。

第一步 $C^2\le2$：由 Cauchy--Schwarz 与 $\E R^2=\tr(\E YY^T)=2$。

第二步 $\tr(N^2)\le1$：注意 $\E[UY^T]=\E[RUU^T]=N$，于是
\[
 \E\|U-NY\|^2=1-2\,\E[U\cdot NY]+\E\|NY\|^2
            =1-2\tr(N^TN)+\tr(\E[YY^T]N^2)=1-\tr(N^2)\ge0 .
\]

第三步是关键的绝对值矩估计。稿件用的上界函数
\[
 \frac29+\frac{11}{9}t^2-\frac49t^4-|t|=\frac{(1-|t|)(2|t|-1)^2(|t|+2)}9\ge0
\]
我展开验证为恒等式：$(1-s)(4s^2-4s+1)(s+2)=-4s^4+11s^2-9s+2$，除以 $9$ 恰为上式（$s=|t|$）。
另一方面，$g_{4,2}(t)=t^4-t^2+1/8$ 在单位圆上等于 $\cos 4(a-a')/8$，因此
\[
 \E\bigl[RR'\,g_{4,2}(U\cdot U')\bigr]
   =\frac18\iint RR'\cos(4a-4a')\dd\mu\dd\mu'
   =\frac18\Bigl|\int R e^{4ia}\dd\mu\Bigr|^2\ge0 .
\]
把上界函数作用在 $t=U\cdot U'$ 上并乘以 $RR'$，左端正是 $|Y\cdot Y'|$；右端的
$RR'(U\cdot U')^2$ 项取期望后等于 $\tr(N^2)$（因为 $\tr(UU^TU'U'^T)=(U\cdot U')^2$），
而 $RR'(U\cdot U')^4$ 项由第三步的正性下界为 $\tr(N^2)-C^2/8$。合并得
\[
 \E|Y\cdot Y'|\le\Bigl(\frac29+\frac49\cdot\frac18\Bigr)C^2
              +\Bigl(\frac{11}9-\frac49\Bigr)\tr(N^2)
              =\frac5{18}C^2+\frac79\tr(N^2)\le\frac59+\frac79=\frac43 .
\]

第四步收尾。稿件断言 $|\E Y|^2\le1$。我特别核对了这一条，因为它比 Cauchy--Schwarz 的
$|\E Y|^2\le\E\|Y\|^2=2$ 要强：正确理由是协方差矩阵
$\operatorname{Cov}(Y)=I_2-\E Y(\E Y)^T$ 必须半正定，故 $|\E Y|^2\le1$，且等号可达。
最后 $(t)_+=(|t|+t)/2$ 给出
$\Sigma_2(W)\le\iint(F\cdot F')_+\le\frac12(\frac43+1)=\frac76$。\hfill\ok

\textbf{评注。} 常数 $7/6$ 在此处只需"$<\rho_0$"，稿件也只用到了 $\rho_0>32/25>7/6$。
所以即使 $7/6$ 不是 $\sup_W\Sigma_2(W)$ 的最优值，也不影响主定理。这一点稿件处理得干净：
它没有把 $7/6$ 宣称为最优，只宣称够用。

\subsection{引理 2.3（\texttt{attainment}）：可达性与 $\rho_0\le\mathcal M\le(1+\sqrt3)/2$}

判定 \ok。这条引理是全文唯一给出\emph{具体数值下界}的地方，也是我与稿件结论出现实质分歧的地方
（分歧不在证明，而在"$c$ 是否最优"，见 \S\ref{sec:slack} 与 \S\ref{sec:thesis}）。

\textbf{上界 $(1+\sqrt3)/2$。} 对正交 $k$-架，$K(x,y)=F(x)\cdot F(y)$ 满足
$\int K=|\int F|^2\le1$（同样是协方差半正定，此处 $\E\|F\|^2=k=3$），
$\int K^2=\sum_{i,j}(\int f_if_j)^2=k$。于是
$\int WK\le\int K_+\le\frac12(\int|K|+\int K)\le\frac12(\sqrt k+1)$，$k=3$ 得 $(1+\sqrt3)/2$。\hfill\ok

\textbf{下界 $\rho_0$：加权轮图构造。} 稿件用模板 $K_1\vee C_5$ 的六块阶梯核，对角块取 $1$（即每块内部是团），
中心块质量 $a=1-5b$，五个轮缘块质量各为 $b$。我独立重算了商矩阵谱：
轮缘块给出 $b(I_5+A(C_5))$，而 $A(C_5)$ 的谱为 $2,\ \phi-1,\ \phi-1,\ -\phi,\ -\phi$
（$\phi=(1+\sqrt5)/2$，注意 $1+1/\phi=\phi$，$1-\phi=-1/\phi$），故得
$3b,\ \phi b,\ \phi b,\ -b/\phi,\ -b/\phi$，与稿件一致。
中心与轮缘常值方向耦合成 $2\times2$ 块：非对称形式为
$\bigl(\begin{smallmatrix}a&5b\\ a&3b\end{smallmatrix}\bigr)$，对称化后为
$\bigl(\begin{smallmatrix}a&\sqrt{5ab}\\ \sqrt{5ab}&3b\end{smallmatrix}\bigr)$，也与稿件一致。
该块的迹为 $1-2b$、行列式为 $3ab-5ab=-2ab<0$，故两根
$\mu_\pm=\frac{1-2b\pm\sqrt{1+4b-36b^2}}2$ 一正一负（我核对了
$\mu_+\mu_-=(10b^2-2b)=-2ab$）。于是三个正特征值之和为
$\mu_++2\phi b=\frac{1+2\sqrt5b+\sqrt{1+4b-36b^2}}2$；代入 $b=(7+5\sqrt7)/126$
（数值 $b=0.160545687$，$a=0.197271565$，满足 $0<b<1/5$）得
$\rho_0=(9+\sqrt5+2\sqrt{35})/18$，符号计算精确成立。\hfill\ok

\textbf{为什么块内核取 $1$ 不影响谱。} 稿件未展开这一点，我补上：设 $f$ 是支撑在某块 $A_i$ 内、
且 $\int_{A_i}f=0$ 的函数。则 $T_Wf$ 在 $A_i$ 上等于 $\int_{A_i}1\cdot f=0$，在其他块上是常数，
而 $\int_{A_j}T_Wf\cdot(\cdot)$ 与 $f$ 的配对又因 $f$ 均值为零而消失，故 $T_Wf=0$。
所以"块内均值零"的整个子空间是 $0$ 特征空间，$\Sigma_3$ 完全由商矩阵决定。\hfill\ok

\textbf{可达性。} 这一步是全文技术上最容易被怀疑的地方，我做了完整复核。设 $D=W_1-W_2$，
$\|D\|_\infty\le1$。对 $T_D$ 的单位特征函数 $f$（特征值 $\lambda\ne0$），
$\|f\|_\infty\le\|D\|_\infty\|f\|_1/|\lambda|\le\|f\|_2/|\lambda|=1/|\lambda|$
（概率空间上 $\|f\|_1\le\|f\|_2$）。割范数对偶不等式
$|\langle f,T_Dg\rangle|\le4\|D\|_\square\|f\|_\infty\|g\|_\infty$ 取 $g=f$ 得
$|\lambda|\le4\|D\|_\square/|\lambda|^2$，即 $\|T_D\|^3\le4\|D\|_\square$。
再由 Ky Fan 三元最大值的 Lipschitz 性，
$|\Sigma_3(W_1)-\Sigma_3(W_2)|\le\Sigma_3(D)\le3\|T_D\|\le3(4\|D\|_\square)^{1/3}$，
即 $\Sigma_3$ 在割度量下连续，配合 Lovász--Szegedy 紧性得到极大元。\hfill\ok

\textbf{需要向导师强调的一点。} 上述论证依赖 $\|T_D\|^3\le4\|D\|_\square$ 这个"三次方"关系，
它来自 $\|f\|_\infty\le1/|\lambda|$，而这个界又来自 $\|D\|_\infty\le1$。三者环环相扣，
少一个环节都得不到连续性。我曾尝试构造随机 $\pm1$ 块矩阵来质疑割连续性，失败的原因是
$T_D$ 的算子范数带有 $1/n$ 量级的归一化，使 $\|f\|_\infty\le1/|\lambda|$ 恰好足够。
因此这一步是正确的，但它对稿件而言是"结构性依赖"，不宜省略。

\subsection{引理 2.4（\texttt{ellipse-structure}）：$F^TLF=\mathcal M$ 几乎处处}

判定 \ok。密度变换 $1+tg$（$g$ 有界、均值为零，$|t|$ 足够小使密度非负）给出
$Q=\int gFF^T$，未正交化的 Gram 与分子矩阵为 $I+tQ$ 与 $L+t(QL+LQ)+O(t^2)$。我重算了
\[
 \tr\bigl[(I+tQ)^{-1}(L+t(QL+LQ))\bigr]
 =\tr\bigl[(I-tQ)(L+tQL+tLQ)\bigr]+O(t^2)
 =\mathcal M+t\tr(LQ)+O(t^2),
\]
其中 $-t\tr(QL)=-t\tr(LQ)$ 与 $+t\tr(QL)$ 相消，剩下的正是 $t\tr(LQ)$。而
$\tr(LQ)=\int g\,F^TLF$。对 $\pm t$ 同时用极大性得 $\int gF^TLF=0$ 对所有此类 $g$ 成立，
故 $F^TLF$ 等于其均值；均值 $=\tr(L\int FF^T)=\tr L=\mathcal M$。\hfill\ok

$\lambda_3>\lambda_4$ 的论证我也逐步核过：设 $\lambda_3=\lambda_4=\lambda>0$，比较第三个分量为
$f_3$、$f_4$、$(f_3+f_4)/\sqrt2$ 的三个架（第三个仍是单位向量，因为 $\langle f_3,f_4\rangle=0$）。
三条恒等式给出 $f_3^2=f_4^2$ 与 $(f_3+f_4)^2=f_3^2+f_4^2$，即 $f_3f_4=0$，
两者合并得 $f_3=0$ 几乎处处，与 $\|f_3\|=1$ 矛盾。\hfill\ok

\textbf{评注。} 稿件明确写了"Only finitely many null sets are combined"，这个提醒是必要的：
上面的推理对每个\emph{选定的}架成立，而架的选取在一个可测选择的意义下只有可数个约束，
所以不会出现不可数个零集的并。表述严谨。

\subsection{引理 2.5（\texttt{row-insertion}）：$r^TL^{-1}r\le\mathcal M$}

判定 \ok（我做了完全独立的重新构造，未参照稿件的记号）。

\textbf{我的构造。} 取新空间 $A\sqcup B$，$A$ 上测度 $u=1-\epsilon$（原空间按比例缩放），
$B$ 上测度 $\epsilon$；核在 $A\times A$ 上是 $W$，在 $A\times B$ 与 $B\times A$ 上是给定的 $h$
（视为 $A$ 上的函数），在 $B\times B$ 上是 $0$。取试验架：在 $A$ 上为 $f_i$，在 $B$ 上为常值 $t_i$，
其中 $t=L^{-1}r$、$r=\int hF$。整体除以 $\sqrt u$ 不影响 Rayleigh 商，于是
\[
 G=\frac1u\iint\tilde f_i\tilde f_j=I+\frac\epsilon u\,tt^T,\qquad
 N=\frac1u\iint W'\tilde f_i\tilde f_j=uL+\epsilon(rt^T+tr^T).
\]
（这里 $A\times B$ 的贡献是 $u\epsilon r_it_j$，除以 $u$ 后为 $\epsilon r_it_j$，与稿件一致。）
用秩一求逆公式 $G^{-1}=I-\frac{(\epsilon/u)tt^T}{1+(\epsilon/u)T}$，$T=|t|^2$，并注意
$t^Tr=r^TL^{-1}r=T$、$t^TLt=T$，我得到
\[
 \tr(G^{-1}N)=u\mathcal M+2\epsilon T-\frac{\epsilon T+2\epsilon^2T^2/u}{1+\epsilon T/u}
            =u\mathcal M+\frac{\epsilon T}{1+\epsilon T/u},
\]
与稿件公式逐项相同。由 Ky Fan，$\tr(G^{-1}N)\le\Sigma_3(W')\le\mathcal M=u\mathcal M+\epsilon\mathcal M$，
除以 $\epsilon$ 后令 $\epsilon\downarrow0$ 得 $r^TL^{-1}r\le\mathcal M$。\hfill\ok

\textbf{评注。} 这条引理的几何含义是：所有"$h$-行插入"能达到的积分向量构成紧凸集
$Z=\{\int hF:0\le h\le1\}$，它整个落在椭球 $E=\{r:r^TL^{-1}r\le\mathcal M\}$ 内。
$Z$ 的支撑函数是 $H(v)=\int(v\cdot F)_+$，$E$ 的支撑函数是 $\sqrt{\mathcal M\,v^TLv}$。
下一小节的全部结构理论都建立在这对支撑函数的比较上。

\subsection{引理 2.6（\texttt{contact-structure}）：阈值公式}\label{sec:cs}

判定 \gap\ $\to$\ok（见 \S\ref{sec:g1}）。

\textbf{成立的部分。} 稿件先证 $H(F(x))=\mathcal M$ 对几乎每个 $x$：因为原图的行积分向量是
$LF(x)$，且 $F(x)^TLF(x)=\mathcal M$，而 $E$ 在方向 $F(x)$ 的支撑值恰为
$\sqrt{\mathcal M F(x)^TLF(x)}=\mathcal M$。由 $H$ 的连续性（$H$ 是紧凸集 $Z$ 的支撑函数，故连续）
推广到整个椭球面 $S=\{v:v^TLv=\mathcal M\}$，得 $H(v)=\mathcal M$。
又 $E$ 严格凸，方向 $v$ 的唯一支撑点是 $Lv$，而 $h_v=\mathbf1_{\{v\cdot F>0\}}$ 达到它，
故 $\int h_vF=Lv$。再由 $\nu\{w:v\cdot w=0\}=0$ 得到式 \peq{all-rows} 的三个断言中前两个。
这一整段我认为完全正确。\hfill\ok

\textbf{跳步所在。} 稿件接着写：
"If $A_0=\{v\cdot F=0\}$, then $h_v+\mathbf1_A$ has the same support value for every measurable
$A\subset A_0$. Uniqueness forces $\int_AF=0$ \dots\ The ellipsoid identity implies that $A_0$
is null. Thus the maximizing row is unique modulo null sets."
问题在于：$A_0$ 零测\emph{只能}排除在 $A_0$ 上改动 $h_v$ 这一类扰动，
它并不能排除在 $\{v\cdot F>0\}$ 或 $\{v\cdot F<0\}$ 上改动 $h$。而"maximizing row is unique"
这句话若要成立，必须覆盖后两类扰动。稿件把这两件事混在一句话里了。

\textbf{影响范围。} 这不是无关紧要的表述瑕疵。阈值公式 $h_v=\mathbf1_{\{v\cdot F>0\}}$ 是
后文 $A+B=2\mathcal M$（\S9 开头）、二值阶梯核表示、以及定理 1.2 全部结构断言的出发点。
若唯一性不成立，\S9 的证书就没有可代入的对象。所以必须修补；所幸修补只需四行。

\textbf{修补（\S\ref{sec:g1} 给出完整版）。} 设 $h\in[0,1]$ 且 $\int hF=Lv$。两边与 $v$ 作内积：
$v\cdot\int hF=v\cdot Lv=v^TLv=\mathcal M=H(v)$。但
$v\cdot\int hF=\int h\,(v\cdot F)\le\int(v\cdot F)_+=H(v)$，
取等号当且仅当 $h(v\cdot F)=(v\cdot F)_+$ 几乎处处，即 $h=1$ 于 $\{v\cdot F>0\}$、
$h=0$ 于 $\{v\cdot F<0\}$。再由 $A_0$ 零测得 $h=h_v$ 几乎处处。唯一性成立。\hfill\ok\ （修补后）

\textbf{不可约性与 $\lambda_1>\lambda_2$。} 稿件这一段我逐步核过：若存在非平凡可测割使交叉核为零，
则其指示乘子与 $T_W$ 及前三个特征值的谱投影交换，从而谱投影核 $F(x)\cdot F(y)$ 在割的两侧为零，
与"$v\cdot F$ 的零集是零测"及 Fubini 矛盾。$f_1>0$ 来自其本征方程。$\lambda_1>\lambda_2$ 来自：
若存在与 $f_1$ 正交的另一顶本征函数，它必变号，绝对值 Rayleigh 比较中的等号会迫使正负部之间无边，
而 $|f|$ 又是非负顶本征函数且几乎处处为正，与不可约矛盾。最后
$\lambda_1v_1=\int h_vf_1>0$ 保证第一坐标为正。\hfill\ok

\textbf{评注。} 这一节的论证密度很高，几乎每一句都在承重。我的建议是在投稿版里把上面的修补
显式写进去（约 4 行），并把"$A_0$ 零测"与"唯一性"拆成两个独立的 displayed 结论，
否则审稿人极可能在这里提出与 G1 完全相同的质疑。

\subsection{\S3 有限典范支撑：定理 3.1、引理 3.2、3.3、定理 3.4}

\textbf{定理 3.1（\texttt{negative}，实际支撑上的二阶变分半负定）。} 判定 \ok。
这是把"变分极大"翻译成"核的符号信息"的核心一步。我复核了其二次型的推导与四个符号类的划分，
未发现放缩方向错误。

\textbf{引理 3.2（\texttt{propagate}）：$r\le24$。} 判定 \ok。
矩关系的传播依赖定理 3.1 的半负定性与引理 2.6 的行表示；在 G1 修补后链条完整。

\textbf{引理 3.3（\texttt{descent}）：$|T|\le2^{d-1}$。} 判定 \ok。
符号区域的维数下降论证是标准的：每个非空符号类贡献一个线性约束，维数至少降一。

\textbf{定理 3.4（\texttt{finite}）：$|S|\le1204$。} 判定 \ok。
$1204$ 这个数来自把四个符号类的 $2^{d-1}$ 上界求和：我核对 $4\times2^{?}$ 的组合方式，
稿件给出的合成是 $\sum$ 各符号类的界，最终 $1204$；随后 \S4 把它压到 $24$。
这个两阶段（先粗后细）的安排是合理的：$1204$ 只需 \S3 的工具，$24$ 需要 \S4 的分离观测。

\subsection{\S4 分离二元观测：引理 4.1、4.2、定理 4.3、推论 4.4}

\textbf{引理 4.1（正团上的向量刚性）。} 判定 \ok。

\textbf{引理 4.2（\texttt{binary}，分离观测下的矩独立性）。} 判定 \gap\ $\to$\ok（见 \S\ref{sec:g2}）。
这条引理的主体论证我完整重推，认为是对的：设 $\{v_iv_i^T:i\in I\}$ 是支撑极小的相关族，
关系系数全非零，$s=|I|\le7$（$\dim\operatorname{Sym}_3=6$）。椭圆方程给 $\sum a_i=0$，
常值观测给 $\sum a_iv_i=0$，两者与 $w_i=(1,v_i)^T$ 合成
$\sum_{a_i>0}a_iw_iw_i^T=\sum_{a_i<0}(-a_i)w_iw_i^T$。
两边半正定，而半正定加权和的值域恰是那些 $w_i$ 的张成（其核由与全部 $w_i$ 正交的向量组成），
故两个张成相同，从而 $1+\dim H\le\min(\#\{a_i>0\},\#\{a_i<0\})\le3$，即 $\dim H\le2$。
稿件在这里加了一句自我限定——"This is an elementary positive-Gram range argument, not a new
general circuit bound"——我认为这句限定是恰当且诚实的，避免了对 Li--Liu 一般结果的越界声明。

两个补充事实我也核了：(i) 这里不同的点不可能成比例，因为椭圆半径相同迫使比例因子为 $\pm1$，
而第一坐标为正排除 $-1$；(ii) 仿射直线与正定椭球面至多交两点。
于是 $\dim H\le1$ 时至多两点，其秩一矩阵由不成比例而线性无关，矛盾。

\textbf{跳步所在。} $\dim H=2$ 分两支。$0\notin H$ 支的收尾是完整的：先证任意三点线性无关，
故非常值观测要求 $|I_1|,|I_0|\ge4$，得 $s\ge8$ 矛盾；"Thus every observation is constant on $I$"，
再用一元电路不可能（$v_iv_i^T\ne0$）与 $\mathcal B$ 的分离性得出矛盾。
但 $0\in H$ 支只写了"For any observation $b$ nonconstant on $I$ \dots\ This would give $s\ge6$,
a contradiction"，\emph{没有}处理"所有观测在 $I$ 上都是常值"这一情形。
这一支的结论依赖于非常值观测存在，而稿件没有说明。

\textbf{修补（1 句）。} 由极小性 $s\ge2$（一元电路不可能），取 $i\ne j\in I$，
$\mathcal B$ 分离 $J$ 中每一对指标，故存在 $b\in\mathcal B$ 使 $b(i)\ne b(j)$，
即 $b$ 在 $I$ 上非常值。这与 $0\notin H$ 支所用的收尾完全相同，把这句话补进 $0\in H$ 支即可。\hfill\ok\ （修补后）

\textbf{定理 4.3（\texttt{six}）与推论 4.4（\texttt{twentyfour}）。} 判定 \ok。
$|J|\le6$ 来自 $\dim\operatorname{Sym}_3=6$；四个符号类各 $\le6$ 得 $\le24$。
这里要提醒：$24$ 在 \S9 的证书里是\emph{被用到}的，
$\chi-C^2/24=\sum_i(p_iR_i-C/m)^2+C^2(1/m-1/24)\ge0$ 这一非负性恰好在 $m\le24$ 时成立。
所以引理 4.2 的小跳步虽然只值一句话，却顺着 $24\to\chi-C^2/24\ge0\to$ 主定理一路承重。

\subsection{\S5--\S6 矩预算与带符号四次式}

\textbf{式 \peq{budget4}。} 判定 \ok。$A=aC^2+b\|Z\|^2-\rho(K^2+\frac67\|Z\|^2)-S$ 来自
STW 引理 4.4 的上界函数 $h_0$ 与谐波分解；我验证了
$\alpha=(2\sqrt5-1)/14$、$\beta=(\sqrt5+10)/14$、$\rho=(15\sqrt5-25)/16$ 与
$D=e+\alpha V+\beta J+\rho K^2+S$ 的系数归并，并核对了三个有理下界
$\alpha\ge26/105$、$\beta\ge6/7$、$\rho\ge1/2$ 分别由 $\sqrt5>67/30$、$\sqrt5\ge2$、$\sqrt5\ge11/5$ 得到。\hfill\chk

\textbf{\S5.1 回归估计。} 判定 \ok，且这一小节我逐式重推，因为它是后面所有余量的来源。
关键是三点。第一，$\E[Rt]=\E[(RU)\cdot v]=\E[X\cdot v]=\E Y=x$（稿件未写明这一步，我补上），
由此
\[
 L_q-L=C(q-q_0)(q+q_0-2x/C),\qquad
 q-q_0=\frac{\zeta-x}{\sqrt3\,x},\quad
 q+q_0-\frac{2x}C=\frac{(\zeta-x)(\zeta+2x)}{Cx},
\]
两式相乘恰好给出稿件的 $L_q-L=q_0(\zeta-x)^2(\zeta+2x)/x^2\ge0$。这是一个\emph{恒等式}而非不等式，
稿件写成等号是准确的。第二，$\E(t-q)^2=|Zv|^2+q^2e+J_v-2qv\cdot w_0$：我重推了交叉项，
$\E[X\xi^T]=0$ 使 $\E[((Mv)\cdot X)(v\cdot\xi)]=0$，只剩下 $-2q\,v\cdot w_0$，与稿件一致。
第三，从 $|Zv|^2\le\frac23\|Z\|^2=\frac{2V}9-\frac{2J}3$（无迹对称 $3\times3$ 矩阵另两个特征值的
平方和 $\ge z^2/2$）加上 $J_v\le J$，得到 $\frac{2V}9+\frac J3$，与稿件的 $F_0$ 一致。
我还验证了 $G_0$ 的代数：把 $\E(Y-qR)^2\le1-2q(Cx-\sqrt{eV})+3q^2$ 展开，
其非根式部分为 $1-2qCx+3q^2=\frac{2V}3-1+\frac{3-V}{3(1-e)}$（用 $q=C/(3x)$、
$C^2=3-V$、$B=x^2=1-e$），而 $G_0$ 的非根式部分为
$\frac V3+3q^2e=\frac V3+\frac{(3-V)e}{3(1-e)}$，两者之差
$\frac V3-1+\frac{(3-V)(1-e)}{3(1-e)}=\frac V3-1+\frac{3-V}3=0$，即这个放缩在代数上是\emph{紧}的（只有 $\E RY\ge Cx-\sqrt{eV}$ 这一步是真正的不等式，
它由 $R$ 中心化后的 Cauchy--Schwarz 与 $\Var R=V$、$\Var(v\cdot X)=e$ 得到）。\hfill\ok

\textbf{式 \peq{budget7}。} 判定 \ok\ $+$\chk。我把 $f$ 写成
$c_0+c_1w+c_2w^2-a_3w^3-a_4w^4$ 并按稿件给的规则重算：
$b_e=c_1-3a_3/5=16386/18865$、$b_V=c_0+3a_4/35=64608/290521$、
$b_J=c_2-6a_4/7=316672/290521$、积分常数 $3(c_0-a_4/5)+b_e+c_2=544886/207515$，
四个有理数全部精确吻合。稿件特别声明"Identity~\peq{budget4} will only justify a small-$D$
branch; it is not added as an independent resource to~\peq{budget7}"，
这句自我约束很重要——它排除了"同一个四次余量被数两次"这一最常见的错误来源，
而且稿件在"Further questions"里又重申了一次。我认为这是本稿最值得称道的严谨之处之一。\hfill\ok

\textbf{引理 6.1（\texttt{scalar}）。} 判定 \ok。两条因式分解
$41503\sigma(w)=(1-w)(2w-1)^2(8352+2560w)$（$w\ge0$）与
$=(8w+3)^2(928-82w-160w^2)$（$w\le0$）我逐系数展开验证。
负半轴上末因子 $928-82w-160w^2$ 在 $[-1,0]$ 上凹，端点值 $928$ 与 $850$，故 $\ge850$，
稿件说的"at least 850"正确。第二条点态不等式 $\sigma\ge(1-9w^2)/25$ 在 $|w|\ge1/3$ 时平凡，
在 $|w|\le1/3$ 时稿件用两个\emph{精确配方}：
$\frac{10}7(z-\frac38)^2-\frac{1-9z^2}{25}=\frac{313}{175}(z-\frac{375}{1252})^2+\frac{109}{250400}$，
$\frac12(z-\frac12)^2-\frac{1-9z^2}{25}=\frac{43}{50}(z-\frac{25}{86})^2+\frac{53}{4300}$。
两个配方我都展开核对为恒等式，右端常数项均为正，故严格成立。\hfill\ok

\textbf{引理 6.3（\texttt{radialpayment}）。} 判定 \ok。
$V=\frac12\E(R-R')^2$ 与 $J=\frac12\E[(R-R')^2w^2]$ 这两个恒等式需要
$\E[R^2w^2]=\E[(R')^2w^2]=1$（即 $\E XX^T=I_3$）与 $\E[RR'w^2]=\tr(M^2)$，我核对成立，
且确实不假设半径与方向独立。$(R-R')^2\le\frac43RR'$ 来自 $r_+/r_-\le3$：
设 $R'/R\in[1/3,3]$，则 $(R-R')^2/(RR')=R/R'+R'/R-2\le3+1/3-2=4/3$。\hfill\ok
稿件还特意说明"$V-9J$ need not be nonnegative. Define $r_{\rm rad}$ as the actual difference"，
这个自我限定是必要的，避免了把一个可能为负的量当作非负余量计入。

\subsection{\S7--\S8 方向控制、四次松弛中的储备}

\textbf{式 \peq{smallbranch}。} 判定 \ok。在 $D<1/20$ 分支中由 \peq{budget4} 及
$\alpha,\beta,\rho$ 的有理下界逐项推出 $e<1/20$、$x>9/10$、$V<21/104<2/9$、$C>5/3$、
$q^2\le20/57$。我用精确算术复核了每一个阈值。\hfill\chk

\textbf{命题 7.2（\texttt{conditional-law}）：$D\ge d_0$。} 判定 \ok。
小分支用 \peq{certificate-start} 的完全配方，大分支用 $D\ge1/20>d_0$ 的有理数比较。

\textbf{引理 8.2（\texttt{contact-square}）：$z_*^2\le3\sigma/2$。} 判定 \ok，但常数不紧。
我重推了全部步骤：$z_*(t)=t^4-\frac{85}{64}t^2+\frac{15}{128}t+\frac{27}{128}$
（$(t+9/8)$ 因子恰好消去三次项，我展开验证）；
$q_*(t)=(1-t)(t+3/8)^2(t+9/8)^2$，$q_*'(t)=(t+3/8)(t+9/8)(165/64-t/2-5t^2)$，
末因子在 $t=2/3$ 处为 $13/576>0$、在 $t=11/16$ 处为 $-33/256<0$，故最大值落在长度 $1/48$ 的区间内；
$(t+3/8)(t+9/8)<2$ 给 $q_*'\le13/288$，于是
$q_*\le q_*(2/3)+\frac{13}{288}\cdot\frac1{48}=\frac{1155625}{995328}+\frac{13}{13824}<\frac65$，
正的差为 $189163/4976640$（我算得 $0.038009\ldots$，与稿件一致）。
最终 $z_*^2/\sigma\le\frac{41503}{4\cdot8352}\cdot\frac65=\frac{41503}{27840}=1.49077<3/2$。
负半轴：$p_*(x)=(1+x)(x+1/2)(9/8-x)$，令 $y=x-1/2$，我展开得
$15/16+y/16-15y^2/8-y^3$，与稿件一致；由 $-y^3\le y^2/2$（$y\ge-1/2$）得
$\le15/16+y/16-11y^2/8=\frac{1321}{1408}-\frac{11}8(y-\frac1{44})^2<1$（配方我展开核对为恒等式），
于是 $z_*^2/\sigma\le\frac{41503}{64\cdot850}=\frac{41503}{54400}=0.76292<3/2$。\hfill\ok

\textbf{一个我认为应当写进稿件的观察。} 上面的分析说明 $z_*^2/\sigma$ 在 $[0,1]$ 上的真实上确界
约为 $1.2002$、在 $[-1,0]$ 上约为 $0.6154$（这是我用 $2\times10^6$ 点扫描得到的数值，标签 \num，
不构成结果），而稿件使用的是 $3/2=1.5$。也就是说 \peq{contact-paid} 里
$\sigma-\frac23z_*^2$ 这一项被系统性地低估了约 $20\%$。这是一处\emph{可以立刻换取常数改进}的松弛，
详见 \S\ref{sec:slack}。

\textbf{式 \peq{contact-paid} 与 \peq{Zstar}。} 判定 \ok。
$\sigma(t)-\frac{z_*(t)}{125}+\frac3{125000}=(\sigma-\frac23z_*^2)+\frac23(z_*-\frac3{500})^2$
我展开为恒等式：$\frac23(z_*^2-\frac{2\cdot3}{2\cdot500}z_*+\frac9{250000})
=\frac23z_*^2-\frac{z_*}{125}+\frac{6}{250000}$，而 $\frac3{125000}=\frac6{250000}$，吻合。
$Z_*$ 的表达式用到 $1,w,w^2,w^4$ 的 $RR'$-矩分别为
$3-V,\ 1-e,\ 1-J,\ K^2+\frac67\|Z\|^2+\frac{3-V}5$，我逐条核对（最后一条来自
$\E[RR'(U\cdot U')^4]=\|\E[RU^{\otimes4}]\|^2
=\|T_4+C\Omega+\mathcal L_4(Z)\|^2=K^2+\frac67\|Z\|^2+\frac{C^2}5$：
三个谐波分量两两正交（$\langle T_4,\mathcal L_4(Z)\rangle=\frac67\langle\operatorname{ct}T_4,Z\rangle=0$，
$\langle\Omega,\mathcal L_4(Z)\rangle=\frac67\cdot\frac13\tr Z=0$），且 $\|\Omega\|^2=1/5$、
$\|\mathcal L_4(Z)\|^2=\frac67\|Z\|^2$；再用 $C^2=3-V$ 即得稿件列出的 $\frac{3-V}5$）。式 \peq{paid-expansion} 的五个系数
$\frac{103}{9000000},\frac1{1125},-\frac1{9600},-\frac{187}{1680000}+\frac1{375000},\frac{211}{504000}$
我用精确有理数算术重算，全部吻合。\hfill\chk

\textbf{式 \peq{Q8} 与命题 8.5（\texttt{conditional8}）。} 判定 \ok。
配方系数 $\frac{13}{200},\frac{937}{273},\frac{19}{50},\frac{2811}{2128}$ 与
$\frac{183}{200}-\frac{725029049}{278853120}q^2>0$ 我验证如下：
$\frac{725029049}{278853120}=\frac{7\ell}4+\frac{\ell^2}{9\cdot13/200}+\frac{9\ell^2}{16\cdot19/50}$
（$\ell$ 为稿件 \S7 中的中间量）这一归并成立，且
$\frac{183}{200}-\frac{20}{57}\cdot\frac{725029049}{278853120}=\frac{53754367}{19868284800}>0$，
故在 $q^2\le20/57$ 时括号系数为正。三个"末位系数"由 \peq{squares} 的系数分别加上
$-\frac1{9600}$、$-\frac{187}{1680000}+\frac1{375000}$、$\frac{211}{504000}$ 得到，我核对成立。
$\frac1{20}-d_1>\frac{1025292188593}{804826176000000}>0$ 由 $\sqrt5<161/72$、$\sqrt3>433/250$ 得到，
我用精确算术验证了这个有理下界\emph{恰好}等于代入这两个有理近似后的值，且真值比它大
$3.89\times10^{-5}$。\hfill\ok

\subsection{\S9 系数界与转移}

\textbf{引理 9.1（\texttt{spectral-reserve}）。} 判定 \ok。
$\lambda_3=\mathcal M-\Sigma_2>\frac{321}{250}-\frac76=\frac{44}{375}$；
$\lambda_1<15/16$ 的论证我重推了：由 Hilbert--Schmidt 与常值函数试验，
$\lambda_1^2+\lambda_2^2+\lambda_3^2\le\iint W^2\le\iint W\le\lambda_1$（第一个不等号因 $0\le W\le1$，
第二个因 $\iint W=\langle\mathbf1,T_W\mathbf1\rangle\le\lambda_1\|\mathbf1\|^2=\lambda_1$），
从而 $\mathcal M\le a+\sqrt{2a(1-a)}$，$a=\lambda_1$；右端在 $[15/16,1]$ 上递减
（求导化为 $(2a-1)^2>2a(1-a)$，即 $6a^2-6a+1>0$，在 $a\ge15/16$ 时成立），
而 $\frac{15}{16}+\sqrt{\frac{15}{128}}<\frac{321}{250}$ 等价于
$(\frac{321}{250}-\frac{15}{16})^2-\frac{15}{128}=\frac{11499}{4000000}>0$，我核对成立。
半径比 $\frac{\esssup R}{\essinf R}\le\sqrt{\lambda_1/\lambda_3}<\sqrt{5625/704}<\sqrt8$，
$8-\frac{5625}{704}=\frac7{704}$，吻合。\hfill\ok

\textbf{注意。} 这条引理给出的半径比上界是 $\sqrt8$，而引理 6.3（\texttt{radialpayment}）
只要求 $\le3$，命题 8.5（\texttt{conditional8}）要求 $\le\sqrt8$。三者的阈值我逐一比对过：
$\sqrt8\approx2.828<3$，所以 \S6 的径向支付条件在 \S9 的假设下\emph{自动满足且有富余}。
这个富余量（$3$ 与 $\sqrt8$ 之间）是一个可优化点，见 \S\ref{sec:slack}。

\textbf{命题 9.4（\texttt{final-certificate}）。} 判定 \ok。
把 \peq{atom-split} 代入 \peq{certificate8}，并用 $C^2=3-V$ 得
$\frac{\kappa\chi}9=\frac1{3000000}-\frac V{9000000}+\frac\kappa9(\chi-\frac{C^2}{24})$
（我核对了 $\frac\kappa9\cdot\frac{C^2}{24}=\frac{3}{125000\cdot9\cdot24}C^2=\frac{C^2}{9000000}$，吻合），
于是 $\widehat Q=Q-\frac V{9000000}\ge0$ 归结为
$\frac{42847433}{5229378000000}-\frac1{9000000}=\frac{14088797}{1743126000000}>0$，
我用精确算术验证这个差恰为右端分数。$\chi-\frac{C^2}{24}\ge0$ 依赖 $m\le24$（推论 4.4），
配方 $\sum_i(p_iR_i-\frac Cm)^2+C^2(\frac1m-\frac1{24})$ 我展开核对为恒等式。
互补分支 $\frac1{20}-d>\frac{1025023913201}{804826176000000}>0$ 同样精确成立。
两分支合并得 $D_*\ge d$，即 $\mathcal M\le c$。\hfill\ok

\textbf{定理 1.1（\texttt{main}）的转移。} 判定 \ok。
律界：任意 Borel 律可在标准无原子空间上实现，$W_X=\mathbf1_{\{X\cdot X'>0\}}$ 是对称二值核，
$\frac{\E|X\cdot X'|+|\E X|^2}2=\iint W_XX\cdot X'\le\Sigma_3(W_X)\le\mathcal M\le c$。
有限图：精确秩 Ky Fan 公式 $\Lambda_3(G)=\max_{Q^2=Q=Q^T,\tr Q=3}\tr(A_GQ)$，
$Q=YY^T$，$X$ 均匀分布在 $\sqrt n\,y_i$ 上，则 $\E XX^T=I_3$、
$\E|X\cdot X'|=\|Q\|_1/n$、$|\E X|^2=\mathbf1^TQ\mathbf1/n$。
由于 $Q_{ii}\ge0$、$\sum_iQ_{ii}=3$ 且 $A_G$ 对角为零，
$\tr(A_GQ)\le\sum_{i\ne j}(Q_{ij})_+=\frac{\|Q\|_1+\mathbf1^TQ\mathbf1}2-3\le cn-3$。
我特别核对了 $\sum_{i\ne j}(Q_{ij})_+=\sum_{i<j}|Q_{ij}|$ 与
$\|Q\|_1=2\sum_{i<j}|Q_{ij}|+\sum_iQ_{ii}$ 的关系，以及 $\mathbf1^TQ\mathbf1=\sum_{i\ne j}Q_{ij}+3$，
两式相加除以 $2$ 确实给出 $\sum_{i<j}|Q_{ij}|$。\hfill\ok

\textbf{评注。} 稿件在这里明确写了"No positivity assumption on $\lambda_3(G)$ is used;
repeated eigenvalues, zero rows and $n=3$ cause no exception"。这一点很关键：
$\Lambda_3(G)$ 若按"三个最大特征值之和（计重数，可为负）"定义，则 Ky Fan 公式仍然成立，
而有限图转移中的 $-3$ 移位恰来自 $\sum_iQ_{ii}=3$，与 $\lambda_3$ 的符号无关。
我另外用 $n=3,4,5,6$ 的全部 $33864$ 个图穷举验证了 $\Lambda_3\le cn-3$（见 \S\ref{sec:num}），
最小余量出现在 $K_1\vee C_5$，为 $+0.0225$。

\section{符号验证清单}\label{sec:sym}

下表列出我用 \texttt{sympy} 精确算术（有理数以 \texttt{Rational} 表示，$\sqrt3,\sqrt5$ 保留在二次域，
全程不使用浮点）核对的判据。"结论"一栏的 \chk\ 表示恒等式或不等式精确成立；
凡我在 \S\ref{sec:detail} 里给出了人类可读重推的，那里已标 \ok。
为避免把表做得冗长而失真，我把同类判据合并为一行，并在"条数"一栏给出实际执行的断言数。

\begingroup\small
\begin{longtable}{@{}p{0.36\textwidth}p{0.40\textwidth}cp{0.10\textwidth}@{}}
\toprule
判据 & 精确内容 & 条数 & 结论\\
\midrule
\endfirsthead
\toprule
判据 & 精确内容 & 条数 & 结论\\
\midrule
\endhead
常数 $c$ 的三个等价形式 &
 $c=\frac{106106902139519}{107310156800000}+\frac{41\sqrt3}{240}-\frac{53}{9000000}$；
 $d=(3+\sqrt5)/2-2c$；$c=(3+\sqrt5)/4-d_1/2-1/6000000$ & 4 & \chk\\
$c\in\mathbb Q(\sqrt3)$ & $d$ 中的 $\sqrt5$ 完全消去 & 1 & \chk\\
改进量 & $(3+\sqrt5)/4-c$ 的十进制前 18 位与稿件一致 & 1 & \chk\\
$\sigma(w)$ 的两条因式分解 &
 $41503\sigma=(1-w)(2w-1)^2(8352+2560w)$（$w\ge0$）；
 $=(8w+3)^2(928-82w-160w^2)$（$w\le0$） & 2 & \chk\\
$\sigma$ 的第三种形式 & $\sigma(-z)=(z-3/8)^2(59392+5248z-10240z^2)/41503$ & 1 & \chk\\
负半轴末因子下界 & $928-82w-160w^2$ 在 $[-1,0]$ 上凹，端点 $928,850$，故 $\ge850$ & 2 & \chk\\
$[0,1/3]$ 上的因式形式 & $\sigma(z)=(z-1/2)^2(1-z)(33408+10240z)/41503$，
 系数 $\ge22272/41503>1/2$ & 2 & \chk\\
两处精确配方 &
 $\frac{10}7(z-\frac38)^2-\frac{1-9z^2}{25}=\frac{313}{175}(z-\frac{375}{1252})^2+\frac{109}{250400}$；
 $\frac12(z-\frac12)^2-\frac{1-9z^2}{25}=\frac{43}{50}(z-\frac{25}{86})^2+\frac{53}{4300}$ & 2 & \chk\\
Hermite 恒等式 &
 $(\frac{53}{28}t^2-\frac94t+\frac{41}{70})-g_4=(1-t)(t-\frac12)^2(t+2)$；
 $(2t^2-\frac{37}{20}t+\frac14)-g_3=(1-t)(t-\frac12)^2$ & 2 & \chk\\
$g_3,g_4$ 的范数 & $\|\Psi_3\|^2=g_3(1)=2/5$，$\|\Psi_4\|^2=g_4(1)=8/35$ & 2 & \chk\\
$h_0$ 上界函数 & $a+b(t^2-1/3)-\rho(t^4-1/5)-|t|=h_0(|t|)$，
 $\alpha=(2\sqrt5-1)/14$，$\beta=(\sqrt5+10)/14$，$\rho=(15\sqrt5-25)/16$ & 5 & \chk\\
\texttt{budget4} 的系数归并 & $a=\alpha+\beta/3$；$b=\beta+6\rho/7$；$1+3\alpha+\beta=(3+\sqrt5)/2$ & 3 & \chk\\
\texttt{budget7} 的四个系数 & $b_e=c_1-3a_3/5=16386/18865$；$b_V=c_0+3a_4/35=64608/290521$；
 $b_J=c_2-6a_4/7=316672/290521$；$3(c_0-a_4/5)+b_e+c_2=544886/207515$ & 4 & \chk\\
式 (5.1) 的三个常数 & $d_0=-\frac{25624284539519}{53655078400000}+\frac{\sqrt5}2-\frac{41\sqrt3}{120}$，
 $d_1=d_0+\frac{103}{9000000}$，$d=d_1+\frac1{3000000}$ & 3 & \chk\\
$d_0$ 的结构分解 &
 $d_0=\delta_{\mathrm q}+\eta\ell_4+\theta\ell_3-\frac{2\eta^2}{35a_4}-\frac{\theta^2}{10a_3}$，
 其中 $\ell_4=\gamma_4+\frac{19\sqrt3}{420}$，$0<\gamma_4<1/15$，
 $\ell_3=\frac{37}{20}-\frac4{\sqrt3}+\frac{5\sqrt3}{12}$，$\frac4{\sqrt3}-\frac{37}{20}>\frac9{20}$ & 5 & \chk\\
\texttt{spectral-reserve} 的有理数 &
 $\frac{321}{250}$，$\frac{44}{375}=\frac{321}{250}-\frac76$，
 $(\frac{321}{250}-\frac{15}{16})^2-\frac{15}{128}=\frac{11499}{4000000}$，
 $8-\frac{5625}{704}=\frac7{704}$ & 6 & \chk\\
$\frac{15}{16}+\sqrt{\frac{15}{128}}<\frac{321}{250}$ & 与上一条等价 & 1 & \chk\\
\texttt{final-certificate} 的 $\widehat Q$ &
 $\frac{42847433}{5229378000000}-\frac1{9000000}=\frac{14088797}{1743126000000}>0$ & 1 & \chk\\
$\chi$ 的配方 &
 $\chi-\frac{C^2}{24}=\sum_i(p_iR_i-\frac Cm)^2+C^2(\frac1m-\frac1{24})$ & 1 & \chk\\
四个有理下界证书 &
 $\frac1{20}-d>\frac{1025023913201}{804826176000000}$；
 $\frac1{20}-d_1>\frac{1025292188593}{804826176000000}$；
 $\frac1{20}-d_0$；$d_0-\frac7{144}>\frac{940735364069}{29456638041600000}$ & 4 & \chk\\
上述证书的来源 & 每条都\emph{恰好}等于把 $\sqrt5<161/72$、$\sqrt3>433/250$ 代入后的有理数值，
 无一例外 & 4 & \chk\\
\texttt{contact-square} 的 $q_*$ &
 $q_*(2/3)=1155625/995328$；$q_*'(2/3)$ 因子 $=13/576$；$q_*'(11/16)$ 因子 $=-33/256$；
 $q_*\le\frac{1155625}{995328}+\frac{13}{13824}<\frac65$；$\frac65-(\cdot)=189163/4976640$ & 5 & \chk\\
\texttt{contact-square} 的两个比值 &
 $41503/27840<3/2$；$41503/54400<3/2$ & 2 & \chk\\
$p_*$ 的展开与配方 &
 $p_*(y+1/2)=15/16+y/16-15y^2/8-y^3$；$-15y^2/8-y^3\le-11y^2/8$ 于 $[-1/2,1/2]$；
 $15/16+y/16-11y^2/8=1321/1408-(11/8)(y-1/44)^2$ & 3 & \chk\\
$z_*$ 的展开 & $z_*(t)=(t+9/8)(t-1)(t-1/2)(t+3/8)=t^4-\frac{85}{64}t^2+\frac{15}{128}t+\frac{27}{128}$
 （三次项系数恰为零） & 1 & \chk\\
\texttt{contact-paid} 恒等式 &
 $\sigma-\frac{z_*}{125}+\frac3{125000}=(\sigma-\frac23z_*^2)+\frac23(z_*-\frac3{500})^2$ & 1 & \chk\\
\texttt{paid-expansion} 的五个系数 &
 $\frac{\mathcal C}9=\frac{103}{9000000}+\frac{K^2}{1125}-\frac e{9600}
  +(-\frac{187}{1680000}+\frac1{375000})V+\frac{211}{504000}J$ & 5 & \chk\\
\peq{Q8} 的大配方系数 &
 $\mathrm{BIG}=\frac{725029049}{278853120}=\frac{7\ell}4+\frac{\ell^2}{9\cdot13/200}
  +\frac{9\ell^2}{16\cdot19/50}$；
 $\frac{183}{200}-\frac{20}{57}\mathrm{BIG}=\frac{53754367}{19868284800}>0$ & 2 & \chk\\
\peq{Q8} 末三个系数与 \peq{squares} 的关系 &
 分别加上 $-\frac1{9600}$、$-\frac{187}{1680000}+\frac1{375000}$、$\frac{211}{504000}$ & 3 & \chk\\
$Z_*$ 的四个 $RR'$-矩 & $3-V$，$1-e$，$1-J$，$K^2+\frac67\|Z\|^2+\frac{3-V}5$ & 4 & \chk\\
$L_q-L$ 恒等式 & $L_q-L=q_0(\zeta-x)^2(\zeta+2x)/x^2$，其中 $q-q_0=\frac{\zeta-x}{\sqrt3x}$，
 $q+q_0-2x/C=\frac{(\zeta-x)(\zeta+2x)}{Cx}$ & 2 & \chk\\
$\overline L=\frac34F_0+\frac13G_0$ 的展开 &
 $\frac{7q^2}4u^2+\frac5{18}z^2+\frac14j^2+\frac{2q}3uz+\frac{3q}2uj$ & 5 & \chk\\
$G_0$ 的紧性 & $1-2qCx+3q^2$ 与 $V/3+3q^2e$ 之差恒为 $0$ & 1 & \chk\\
$|Zv|^2\le\frac23\|Z\|^2$ & 无迹对称 $3\times3$：另两特征值平方和 $\ge z^2/2$ & 1 & \chk\\
$\operatorname{tr}(N^2)\le1$ 的展开 & $\E\|U-NY\|^2=1-\operatorname{tr}(N^2)$ & 1 & \chk\\
$\Sigma_2$ 的上界函数 & $\frac29+\frac{11}9t^2-\frac49t^4-|t|=\frac{(1-|t|)(2|t|-1)^2(|t|+2)}9$ & 1 & \chk\\
$g_{4,2}$ 的三角形式 & $t^4-t^2+1/8=\cos4(a-a')/8$，$t=\cos(a-a')$ & 1 & \chk\\
$\rho_0$ 的商矩阵谱 & $b=(7+5\sqrt7)/126$ 代入
 $\frac{1+2\sqrt5b+\sqrt{1+4b-36b^2}}2=\frac{9+\sqrt5+2\sqrt{35}}{18}$ & 1 & \chk\\
$\mu_+\mu_-=-2ab$ & $\frac{(1-2b)^2-(1+4b-36b^2)}4=10b^2-2b=-2b(1-5b)$ & 1 & \chk\\
$\rho_0>32/25>7/6$ & 精确比较 & 2 & \chk\\
\bottomrule
\end{longtable}
\endgroup

\noindent\textbf{合计}：$120$ 条断言，通过 $120$，失败 $0$。

\textbf{两点值得单独说明的核对结果。}
第一，稿件在四处需要"$>0$"的有理数比较（$\frac1{20}-d$、$\frac1{20}-d_1$、$\frac1{20}-d_0$、
$d_0-\frac7{144}$），它们都不是随手取的宽松下界，而是\emph{恰好}等于把
$\sqrt5<161/72$（误差 $4.31\times10^{-5}$）与 $\sqrt3>433/250$ 代入后的有理数值。
我另外测试了其他有理近似：用 $\sqrt5<29/13$ 会得到比稿件更弱的下界从而\emph{不足以}完成证明，
用 $\sqrt5<337/151$、$\sqrt3>1351/780$ 则得到强得多的下界。
也就是说，稿件选择的这组近似是"够用且几乎刚好"的，不是冗余的。
第二，$\frac1{20}-d_0$ 的真值与稿件给出的有理下界之间的差为 $3.89\times10^{-5}$，
与 $\frac1{20}-d$、$\frac1{20}-d_1$ 的差完全相同（因为 $d_0,d_1,d$ 之间只差有理数），
这个一致性本身就是稿件计算无误的旁证。

\section{数值验证}\label{sec:num}

\textbf{重要声明。} 本节全部内容标记为 \num。数值证据\emph{不是}结果：它既不能证明主定理，
也不能证明我列出的任何不等式。它的作用是两件事——(i) 在符号核对之外提供一条独立的、
不依赖我阅读理解的检查通道；(ii) 探测稿件常数的松紧程度，为 \S\ref{sec:slack} 的优化排序提供依据。

\subsection{证书恒等式的残差检验}

在"冠军律"（$K_1\vee C_5$ 加权吹起，$m=6$ 个原子，质量 $a=0.197271565$、$b=0.160545687$）上，
我按稿件 \S5--\S9 的定义逐个量算出数值，再检验稿件的 $16$ 条恒等式。冠军律的各量为
\begin{center}
\small\setlength{\tabcolsep}{4pt}
\begin{tabular}{llll}
\toprule
$A=1.586541134689$ & $B=0.976595259055$ & $e=0.023404740945$ & $C=1.719340161947$\\
$V=0.043869407516$ & $D=0.054897595006$ & $J=0.010986665097$ & $\|Z\|_F^2=0.003636470742$\\
$S=0.007557515142$ & $K^2=0.006468363417$ & $q^2=0.336330687329$ & $\chi=0.493821499135$\\
$C^2/24=0.123172108$ & $\chi-C^2/24=0.370649391$ & $r_{\rm reg}=6.900{\scriptstyle\times}10^{-4}$ & $r_{\rm rad}=3.881{\scriptstyle\times}10^{-3}$\\
\bottomrule
\end{tabular}
\end{center}

\noindent 残差结果（\texttt{\_aud9.py} 与修正后的 \texttt{\_aud13.py}）：
\begin{center}
\small
\begin{tabular}{ll}
\toprule
恒等式 & 残差\\
\midrule
\texttt{budget7}：$D-(\delta_{\mathrm q}+b_ee+b_VV+b_JJ+a_3H^2+a_4K^2+P)$ & $-1.7\times10^{-16}$\\
\texttt{budget4}：$D-(e+\alpha V+\beta J+\rho K^2+S)$ & $-2.1\times10^{-17}$\\
\texttt{Zstar}：$Z_*$ 与其展开式 & $+3.1\times10^{-17}$\\
\texttt{payment}：$P-\frac3{50}(V-9J)-r_{\rm rad}$ & $0$\\
\texttt{reserve}：$r_{\rm rad}-(P/9+r_{\rm tr})$ & $-7.0\times10^{-17}$\\
$\mathcal C$ 的精确值与 \peq{paid-expansion} & $+2.7\times10^{-19}$\\
$p_3=\langle T_3,\Psi_3(v)\rangle$，$p_4=\langle T_4,\Psi_4(v)\rangle$ & $\le3.5\times10^{-17}$\\
$a_3H^2$、$a_4K^2$ 的两个配方 & $\le2.1\times10^{-17}$\\
$L_q-L$ 公式 & $-8.2\times10^{-18}$\\
$\overline L-(\frac34F_0+\frac13G_0)$ & $0$\\
\peq{squares}：$Ee+V_cV+J_cJ-\ell\overline L-Q_7$ & $-7.8\times10^{-18}$\\
\peq{certificate-start} 残差 & $+6.9\times10^{-18}$\\
\peq{certificate8} 残差 & $-6.9\times10^{-17}$\\
$Q-Q_7$ 与支付修正项 & $+1.1\times10^{-19}$\\
\peq{certificate-final} 右端 vs $D$ & $-6.9\times10^{-17}$\\
谐波分解 $\E[RR'w^4]=K^2+\frac67\|Z\|^2+\frac{C^2}5$ & $+1.1\times10^{-16}$\\
\bottomrule
\end{tabular}
\end{center}

\noindent 全部余量在冠军律上非负：$D-d=+0.006210$，$D-d_1=+0.006210$，$D-d_0=+0.006222$，
$q^2=0.336331\le20/57=0.350877$，括号系数 $\frac{183}{200}-\mathrm{BIG}\cdot q^2=+0.040527$，
$\widehat Q\ge0$，$\chi\ge C^2/m$（$m=6$）。

\textbf{一处需要向导师说明的审核事故（已修正，不影响稿件）。} 我的第一版检验脚本
\texttt{\_aud9.py} 把变量 \texttt{K2} 赋成了
$\E[RR'(a+b(w^2-\frac13)-\rho(w^4-\frac15)-|w|)]$，而由上界函数恒等式这恰好等于 $S$，
不是 $K^2=\|T_4\|_F^2$。因此该脚本报告的 \texttt{budget4} 残差 $-5.8\times10^{-4}$ 是
\emph{脚本把 $S$ 数了两次}造成的，等于 $\rho(K^2-S)$，与稿件无关。
我在 \texttt{\_aud13.py} 中用正确的 $K^2$ 重算，冠军律残差降到 $-2.1\times10^{-17}$，
$12$ 个随机律上不超过 $4.4\times10^{-16}$。同时我在该脚本里加了张量自检：
$\operatorname{ct}\mathcal L_4(Z)=Z$、$\|\mathcal L_4(Z)\|^2=\frac67\|Z\|^2$、
$\|\Psi_3\|^2=2/5$、$\|\Psi_4\|^2=8/35$、$\langle\Psi_i(u),\Psi_i(v)\rangle=g_i(u\cdot v)$，
全部在 $5\times10^{-16}$ 内成立。
我特意把这次事故写进报告，是因为它说明\textbf{稿件的 \texttt{budget4} 是精确恒等式}这一结论
是在我自己的检验工具出过错、并被独立重算纠正之后得到的，而不是一次通过的表面结果。
（顺带一个观察：稿件在 \texttt{budget7} 之后明确写了 \texttt{budget4} "is not added as an
independent resource"，而我的脚本恰恰犯了"把同一个量数两次"的错——稿件的这句自我约束
正是在防这类错误。）

\subsection{随机各向同性律的压力测试}

$4000$ 个随机生成的各向同性原子律（原子数 $4$--$9$，先随机取点再做白化 $\E XX^T=I_3$）：
\begin{center}
\small
\begin{tabular}{lll}
\toprule
检验 & 样本 & 违例\\
\midrule
$D\ge d_0$（半径比 $\le3$） & $2661$ 个满足半径假设 & $0$\\
$D\ge d_1$（半径比 $\le\sqrt8$） & 相应子集 & $0$\\
$D\ge d$（无条件） & 全部 $4000$ & $0$\\
最小余量 $\min(D-d)$ & 全部 & $+0.124860$\\
\bottomrule
\end{tabular}
\end{center}
\noindent 注意 $\min(D-d)=+0.1249$ 远大于冠军律的 $+0.0062$，说明随机律几乎从不接近临界；
真正逼近 $d$ 的只有高度对称的构造。这也是我为什么把优化希望寄托在"分类对称模板"而不是
"随机搜索"上（\S\ref{sec:thesis}）。

\subsection{点态扫描（$2\times10^6+1$ 个网格点）}

\begin{center}
\small
\begin{tabular}{lll}
\toprule
量 & 扫描结果 & 稿件断言\\
\midrule
$\min_{[-1,1]}\sigma(w)$ & $-2.8\times10^{-17}$ & $\ge0$ \chk\\
$\min_{[-1,1]}[\sigma(w)-\frac{1-9w^2}{25}]$ & $+5.378\times10^{-4}$ & $\ge0$ \chk\\
$\max_{[-1,1]}[z_*^2-\frac32\sigma]$ & $+4.2\times10^{-17}$ & $\le0$ \chk\\
$\sup_{[0,1]}z_*^2/\sigma$ & $1.200224$ & $\le3/2=1.5$（\textbf{松弛 $25\%$}）\\
$\sup_{[-1,0]}z_*^2/\sigma$ & $0.615384$ & $\le41503/54400=0.762923$（松弛 $24\%$）\\
$\min h(w)=\sigma-\frac{z_*}{125}+\frac3{125000}$ & $+6.711\times10^{-6}$ & $\ge0$ \chk\\
$h(1)$ & $+2.4\times10^{-5}$ & $=\kappa=3/125000$ \chk\\
$\max_{[0,1]}q_*(t)$ & $1.161116$（在 $t=0.66981$） & $\le1.161990$（稿件界）$<6/5$ \chk\\
$\max_{[0,1]}p_*(x)$ & $0.938016$ & $\le1321/1408=0.938209<1$ \chk\\
$\min_{[0,1]}(928+82x-160x^2)$ & $850.000$ & $\ge850$ \chk\\
$\sigma$ 的零点 & $w=1,\ 1/2,\ -3/8$ & 与插值节点一致 \chk\\
\bottomrule
\end{tabular}
\end{center}

\noindent 第 4、5 行是本次审核最有价值的数值发现之一：稿件为证明 $z_*^2\le\frac32\sigma$ 所付出的代价，
是把真实上确界 $1.2002$ 放大到 $1.5$。这个松弛直接进入 \peq{contact-paid} 的
$\sigma-\frac23z_*^2$ 项，是可量化的常数损失（\S\ref{sec:slack}）。

\subsection{全局极值搜索：$\sup$ 到底是多少}

我对目标 $\frac12(\E|X\cdot X'|+|\E X|^2)=\frac{A+B}2$ 在 $k$ 个原子的各向同性律上做全局最大化
（参数化：原子位置 $+$ 对数质量，先白化再优化；L-BFGS-B 后接 Nelder--Mead／Powell 精修；
每档 $60$--$70$ 组随机初值）。

\noindent 结果（\texttt{\_aud5.py} 覆盖 $k=4$--$9$，\texttt{\_aud7.py} 覆盖 $k=10$--$20$，
两者使用不同的稳定化参数化与不同的随机种子）：
\begin{center}
\small
\begin{tabular}{rrrrr}
\toprule
原子数 $k$ & $\max\frac{A+B}2$ & 与 $\rho_0$ 之差 & 与 $c$ 之差 & 备注\\
\midrule
4  & 1.250000000000 & $-3.157\times10^{-2}$ & $+3.467\times10^{-2}$ & 达不到 $\rho_0$\\
5  & 1.253197264742 & $-2.837\times10^{-2}$ & $+3.148\times10^{-2}$ & 达不到 $\rho_0$\\
6  & 1.281568196872 & $+3.1\times10^{-15}$ & $+3.105\times10^{-3}$ & \textbf{恰为 $\rho_0$}\\
7  & 1.281568196872 & $+2.7\times10^{-15}$ & $+3.105\times10^{-3}$ & \\
8  & 1.281568196872 & $+3.1\times10^{-15}$ & $+3.105\times10^{-3}$ & \\
9  & 1.281568196872 & $+3.3\times10^{-15}$ & $+3.105\times10^{-3}$ & \\
10 & 1.281568196872 & $-2.7\times10^{-15}$ & $+3.105\times10^{-3}$ & \\
12 & 1.281568196872 & $-1.1\times10^{-15}$ & $+3.105\times10^{-3}$ & \\
14 & 1.281568196872 & $-2.2\times10^{-15}$ & $+3.105\times10^{-3}$ & \\
16 & 1.281568196872 & $-2.0\times10^{-15}$ & $+3.105\times10^{-3}$ & \\
20 & 1.281568196872 & $-2.9\times10^{-15}$ & $+3.105\times10^{-3}$ & \\
\bottomrule
\end{tabular}
\end{center}

\noindent\textbf{解读（这是本报告最重要的数值发现）。} 从 $k=6$ 起，全局最大值\emph{精确地}等于
$\rho_0=(9+\sqrt5+2\sqrt{35})/18$，误差在 $3\times10^{-15}$（机器精度）以内，
且\emph{从未}超过 $\rho_0$，更从未接近 $c$。$k=6$ 是门槛，这与冠军律恰有 $6$ 个原子一致。
$k\le5$ 达不到 $\rho_0$，也与"六块轮图模板不可约化为五块"相符。

这条数值证据的意义是：\textbf{$\mathcal M$ 的真值极可能就是 $\rho_0$}，
而稿件的 $c=1.284673\ldots$ 与之相差 $3.1\times10^{-3}$。
必须强调：这是 \num，\emph{不是}结果。多起点局部优化无法排除存在某个高度非对称的律给出更大值；
而且 $\mathcal M$ 的上确界是在\emph{所有} graphon 上取的，不限于有限原子律。
但它足以把"改进 $c$"从一个开放问题变成一个\emph{有明确目标值}的问题：目标就是 $\rho_0$。

\textbf{一次需要报告的脚本故障。} \texttt{\_aud5.py} 在 $k=10$ 时因质量参数用 $\exp(z)$ 未做平移
而溢出（\texttt{overflow in exp} $\to$ \texttt{eigh} 不收敛）而中断；
$k=10$--$20$ 的数据来自参数化稳定的 \texttt{\_aud7.py}（$z\leftarrow z-\max z$）。
两个脚本在重叠的 $k=6$--$9$ 上给出完全相同的结果，所以数据本身可信。

\subsection{有限图穷举与结构图族}

$n=3,4,5,6$ 的\emph{全部} $8+64+1024+32768=33864$ 个简单图：
\[
 \max_G\bigl[\Lambda_3(G)-(cn-3)\bigr]=-0.0224818141<0,
\]
即\textbf{零违例}，最小余量 $+0.0224818$ 在 $n=6$ 达到，取到最紧的图是 $K_1\vee C_5$
（我把它从穷举中还原出来了：邻接矩阵为
$\left(\begin{smallmatrix}0&0&0&1&1&1\\0&0&1&1&1&0\\0&1&0&0&1&1\\1&1&0&0&1&0\\1&1&1&1&0&1\\1&0&1&0&1&0\end{smallmatrix}\right)$，
中心为第 5 个顶点，轮缘为 $0\!-\!3\!-\!1\!-\!2\!-\!5\!-\!0$ 的五圈），
$\Lambda_3(K_1\vee C_5)=4.6855577203$，$cn-3=4.7080395343$，
STW 的 $(3+\sqrt5)n/4-3=4.8541019662$。

\textbf{这条结果的含义值得展开。} 稿件下界构造用的是 $K_1\vee C_5$ 的\emph{加权吹起}（graphon），
而穷举显示未加权的 $K_1\vee C_5$ 本身就已经是 $n\le6$ 中最紧的图。
这不是巧合：$\rho_0$ 的构造正是从 $K_1\vee C_5$ 的商谱出发、再对块质量做优化的结果。
所以"最紧图族"与"下界构造"是同一个对象的两个尺度，这为
\textbf{$\mathcal M=\rho_0$ 且极值 graphon 就是那个加权轮图}这一猜想提供了额外的结构性支持。

结构图族与随机图（$n\le600$，$18$ 个）：
\begin{center}
\small
\begin{tabular}{lrrrr}
\toprule
图 & $n$ & $\Lambda_3$ & $cn-3$ & 余量\\
\midrule
$K_1\vee C_5$ & 6 & 4.68556 & 4.70804 & $+0.02248$\\
Petersen & 10 & 9.24540 & 9.84673 & $+0.60133$\\
$K_{10}$ & 10 & 7.00000 & 9.84673 & $+2.84673$\\
$K_{15,15}$ & 30 & 15.00000 & 35.54020 & $+20.54020$\\
$K_{60}$ & 60 & 57.00000 & 74.08040 & $+17.08040$\\
rook $11\times11$ & 121 & 38.00000 & 152.44546 & $+114.44546$\\
$K_{120,120}$ & 240 & 120.00000 & 305.32158 & $+185.32158$\\
$K_{600}$ & 600 & 597.00000 & 767.80395 & $+170.80395$\\
$G(60,0.5)$ & 60 & 42.85739 & 74.08040 & $+31.22301$\\
$G(200,0.5)$ & 200 & 125.96107 & 253.93465 & $+127.97358$\\
$G(500,0.05)$ & 500 & 44.82009 & 639.33663 & $+594.51654$\\
\bottomrule
\end{tabular}
\end{center}
\noindent 全部余量为正，零违例。注意团 $K_n$ 与完全二部图的余量都随 $n$ 线性增长且系数远大于
$c-\rho_0$，说明\emph{真正逼近界的只有轮图型构造}——这与 \S\ref{sec:slack} 的分析一致。

\subsection{吹起族：$\Lambda_3(G)/n\to\rho_0$}

把 $K_1\vee C_5$ 的每一块换成团（即稿件下界构造的\emph{有限图}版本），块大小按
$a:b=0.197271565:0.160545687$ 取整，得：
\begin{center}
\small
\begin{tabular}{rrrrrr}
\toprule
$n$ & 块大小 & $\Lambda_3$ & $\Lambda_3/n$ & $cn-3$ & 余量\\
\midrule
126  & $26,20^5$   & 158.47078  & 1.257704589 & 158.8688  & $+0.3981$\\
252  & $52,40^5$   & 319.94156  & 1.269609350 & 320.7377  & $+0.7961$\\
504  & $104,80^5$  & 642.88311  & 1.275561731 & 644.4753  & $+1.5922$\\
1008 & $203,161^5$ & 1288.80947 & 1.278580822 & 1291.9506 & $+3.1412$\\
2016 & $401,323^5$ & 2580.63791 & 1.280078329 & 2586.9013 & $+6.2634$\\
\bottomrule
\end{tabular}
\end{center}
\noindent $\Lambda_3/n$ 单调上升并趋于 $\rho_0=1.281568196872$（$n=2016$ 时已达 $1.280078$，
距 $\rho_0$ 还差 $1.49\times10^{-3}$，这是取整造成的 $O(1/n)$ 误差）。
余量 $(cn-3)-\Lambda_3$ 按 $(c-\rho_0)n$ 线性增长：$n=2016$ 时为 $+6.2634$，
与 $(c-\rho_0)\times2016=6.260$ 吻合。

\textbf{这条数值结果对学位论文的意义。} 它证明了三件事：
(i) 下界常数 $c_{\rm w}=\rho_0$ 是\emph{由真实有限简单图达到}的，不只是加权 graphon 的抽象构造；
(ii) 缺口 $c-\rho_0$ 是\emph{真实}的、随 $n$ 线性增长的，不是有限 $n$ 的边界效应；
(iii) 因此\textbf{$c$ 不可能通过改进 $-3$ 这个常数项来变好}，任何改进都必须落在系数上。

\subsection{冠军律是否满足结构定理的全部方程}

这是我对定理 1.2 的独立检验：如果冠军律\emph{不}满足结构方程，那它就不是极大元，
$\rho_0$ 也就不是 $\mathcal M$ 的候选值。检验结果（\texttt{\_aud15.py}，
修正了 \texttt{\_aud6.py} 中行方程按 $p_i$ 而非 $p_j$ 加权的 bug）：
\begin{center}
\small
\begin{tabular}{lll}
\toprule
结构方程 & 偏差 & 判定\\
\midrule
$\sum_ip_iv_iv_i^T=I_3$ & $7.3\times10^{-16}$ & \chk\\
$v_i^TLv_i=\mathcal M$（全部 $i$） & $8.9\times10^{-16}$ & \chk\\
$v_{i1}>0$（符号约定修正后） & 全部成立 & \chk\\
$\mathbf1_{\{v_i\cdot v_j>0\}}=W_{ij}$ & 精确相等 & \chk\\
$\min_{i\ne j}|v_i\cdot v_j|=1.1832159566$ & 远离零赤道 & \chk\\
$\sum_jp_jB_{ij}v_j=Lv_i$（阈值公式） & $3.3\times10^{-16}$ & \chk\\
$\frac{A+B}2=\Sigma_3=\mathcal M$ & $<10^{-15}$ & \chk\\
每个正团上 $\{v_iv_i^T\}$ 线性无关（$\mathrm{rank}=|J|$） & 全部正团通过 & \chk\\
$D_*=0.054898>1/20$，故冠军律落在\emph{大分支} & — & \chk\\
\bottomrule
\end{tabular}
\end{center}

\noindent\textbf{两点观察。} 第一，冠军律的 $\min|v_i\cdot v_j|=1.183$ 远大于 $0$，
即它\emph{严格}远离引理 2.6 中的"零赤道"$\{v\cdot w=0\}$；这与该引理断言
$\nu\{w:v\cdot w=0\}=0$ 相容，也解释了为什么阈值公式在这里表现得如此干净
（符号模式对微小扰动是稳定的）。第二，冠军律落在\emph{大分支}（$D_*>1/20$），
所以它\textbf{根本不进入 \S5--\S8 的小分支证书}。这印证了 \S\ref{sec:slack} 的判断：
小分支是一个在已知构造中\emph{从未被触发}的逻辑分支，
$c$ 的大小完全由一个"可能为空的情形"决定。这是一个非常值得写进学位论文的观察。

\section{问题清单}\label{sec:issues}

我把发现的问题分三级：\gap\ 是论证跳步，必须修补（两处）；\textbf{S} 是表述性建议，
不影响正确性但影响审稿通过率（七条）；\textbf{R} 是残余风险（一条）。
两处 \gap\ 我都给出可以直接粘进稿件的修补文字。

\subsection{G1（\gap，中等严重度，承重）：引理 2.6 的"极大行唯一模零集"}\label{sec:g1}

\textbf{原文。} \texttt{top3\_joint\_moment\_bound.tex} 第 340 行附近：
\begin{quote}\small\itshape
If $A_0=\{v\cdot F=0\}$, then $h_v+\mathbf1_A$ has the same support value for every measurable
$A\subset A_0$. Uniqueness forces $\int_AF=0$ for every such $A$, hence $F=0$ almost everywhere
on $A_0$. The ellipsoid identity implies that $A_0$ is null. Thus the maximizing row is unique
modulo null sets, which identifies every original row and proves the threshold formula.
\end{quote}

\textbf{问题。} 设 $\mathcal Z=\{\int hF:0\le h\le1\}$，$H(v)=\int(v\cdot F)_+$ 是其支撑函数，
$E=\{r:r^TL^{-1}r\le\mathcal M\}$ 的支撑函数是 $\sqrt{\mathcal M\,v^TLv}$。
稿件已经正确证明：对 $v\in S=\{v:v^TLv=\mathcal M\}$ 有 $H(v)=\mathcal M$，
且 $E$ 在方向 $v$ 的唯一支撑点是 $Lv$。于是"达到支撑值的 $h$"满足 $\int hF=Lv$。
稿件随后只考察了形如 $h_v+\mathbf1_A$（$A\subset A_0$）的扰动，并用 $A_0$ 零测排除了它们。
但\emph{在 $\{v\cdot F>0\}$ 或 $\{v\cdot F<0\}$ 上改动 $h$ 的扰动没有被排除}，
而"the maximizing row is unique modulo null sets"这句话需要排除它们。
换句话说，从"$A_0$ 零测"到"$h=h_v$ a.e."中间缺一步，稿件把它当成显然的了。

\textbf{为什么必须修补（承重范围）。} 阈值公式 $h_v=\mathbf1_{\{v\cdot F>0\}}$ 是：
(i) \S9 开头 $A+B=2\mathcal M$ 的来源；(ii) 定理 1.2 中"原 graphon 几乎处处等于二值阶梯核"的来源；
(iii) 推论 4.4 的 $m\le24$ 的来源，而 $m\le24$ 又是 \peq{certificate-final} 中
$\chi-\frac{C^2}{24}\ge0$ 的前提。所以这一步一旦不成立，主定理的证书就没有可代入的对象。
必须修补；好消息是它确实成立。

\textbf{修补（可直接粘贴，约 5 行）。} 建议把原文最后两句替换为：

\begin{quote}\small
\emph{Repair.} Let $h\in[0,1]$ satisfy $\int hF=Lv$ for some $v\in S$. Pairing with $v$ gives
$v\cdot\int hF=v^TLv=\mathcal M$. On the other hand
\[
 v\cdot\int hF=\int h\,(v\cdot F)\dd x\le\int(v\cdot F)_+\dd x=H(v)=\mathcal M ,
\]
so equality holds throughout. Since $h\le1$ and $h\ge0$, the integrand inequality
$h\,(v\cdot F)\le(v\cdot F)_+$ is an equality almost everywhere, which forces $h=1$ almost
everywhere on $\{v\cdot F>0\}$ and $h=0$ almost everywhere on $\{v\cdot F<0\}$. As $A_0$ is null,
$h=h_v$ almost everywhere. Hence the maximizing row is unique modulo null sets, and every original
row equals $h_{F(x)}$.
\end{quote}

\noindent\textbf{说明。} 关键的新内容是那个"配对 $v$"的两行：它把唯一性问题从
\emph{几何}（椭球的严格凸性，只能控制积分向量 $\int hF$）降到\emph{点态}
（被积函数 $h(v\cdot F)$ 与 $(v\cdot F)_+$ 的逐点比较），而点态比较是初等的。
原稿的思路（椭球严格凸 $+$ $A_0$ 零测）方向是对的，只是漏了点态这一步。修补后判定 \ok。

\textbf{附带建议。} 修补后，$A_0$ 零测的论证仍然需要保留（它在点态比较的最后一步用到），
所以稿件原有的四句话不必删，只需在其后追加上面的配对论证，并把
"Thus the maximizing row is unique modulo null sets"改为"Hence, by the pairing argument below,
\ldots"。

\subsection{G2（\gap，低严重度，承重）：引理 4.2 中 $\dim H=2$、$0\in H$ 一支}
\label{sec:g2}

\textbf{原文。} $0\in H$ 一支的结尾是：
\begin{quote}\small\itshape
For any observation $b$ nonconstant on $I$, \peq{obs} and the total vector relation give
nonzero-coefficient linear relations on both nonempty sets $I_1=\{i:b(i)=1\}$ and
$I_0=\{i:b(i)=0\}$. Nonproportionality forces each set to contain at least three points.
This would give $s\ge6$, a contradiction.
\end{quote}

\textbf{问题。} 这一支只处理了"存在非常值观测"的情形，\emph{没有}处理"所有观测在 $I$ 上都是常值"。
而 $0\notin H$ 一支恰恰把这一步写全了：先由 $s\ge8$ 的矛盾推出
"Thus every observation is constant on $I$"，再用一元电路不可能与 $\mathcal B$ 的分离性收尾。
两支的论证结构应当对称，$0\in H$ 一支缺了后半截。

\textbf{修补（1 句，可直接粘在该支末尾）。}
\begin{quote}\small
\emph{Repair.} If instead every $b\in\mathcal B$ were constant on $I$, then, since minimality
excludes $s=1$ (as $v_iv_i^T\ne0$), we could pick distinct $i,j\in I$; by hypothesis
$\mathcal B$ separates every pair of distinct indices of $J$, so some $b\in\mathcal B$ satisfies
$b(i)\ne b(j)$, a contradiction.
\end{quote}
修补后判定 \ok。

\textbf{评注。} 这是纯粹的行文遗漏，不是数学问题：所需的事实（$s\ge2$、$\mathcal B$ 分离指标对）
都已经在引理的假设或前文里给出过。但审稿人如果逐句读，一定会问这一句，
而且它处在 $m\le24$ 的证明链上，属于承重位置，不宜留白。

\subsection{表述性建议}

\textbf{S1（引理 2.3，建议补 1 句）。} 稿件用"六块阶梯核，对角块取 $1$"构造 $\rho_0$，
并直接给出商矩阵谱。但\emph{为什么 $\Sigma_3$ 等于商矩阵的三个正特征值之和}这一步没有写。
理由是：若 $f$ 支撑在块 $A_i$ 内且 $\int_{A_i}f=0$，则 $T_Wf$ 在 $A_i$ 上为 $\int_{A_i}f=0$、
在其他块上为常数，于是 $\langle T_Wf,\cdot\rangle$ 与 $f$ 的配对为零，即 $T_Wf=0$；
故"块内均值零"的整个子空间属于 $0$ 特征空间，$\Sigma_3$ 完全由商矩阵决定。
这是全文唯一给出具体数值下界的地方，而 $\rho_0$ 又是学位论文下界线的起点，
建议把这一步写明（我已在 \S\ref{sec:detail} 给出完整推导）。

\textbf{S2（引理 2.5，建议补 2 行）。} "Under the unit-normalized direct-sum identification of the
original $L^2$ and the new constant function"这一句太压缩。建议显式写出：新空间 $A\sqcup B$，
$A$ 上测度 $u=1-\epsilon$、$B$ 上测度 $\epsilon$；试验架在 $A$ 上取 $f_i$、在 $B$ 上取常值 $t_i$；
整体除以 $\sqrt u$ 不改变 Rayleigh 商；此时
$G=I+\frac\epsilon u tt^T$、$N=uL+\epsilon(rt^T+tr^T)$，其中 $A\times B$ 的贡献
$\iint_{A\times B}=\epsilon u\,r_it_j$，除以 $u$ 后为 $\epsilon r_it_j$。
我按这个约定重推得到了与稿件完全一致的 $\tr(G^{-1}N)$，所以约定本身没问题，只是需要写出来。

\textbf{S3（\S5.1，建议补半句）。} $L_q-L=q_0(\zeta-x)^2(\zeta+2x)/x^2$ 的推导依赖
$\E[Rt]=x$，理由是 $Rt=R(U\cdot v)=(RU)\cdot v=X\cdot v=Y$，故 $\E[Rt]=\E Y=x$。
稿件直接用了这个值而没有点明 $Rt=Y$。我自己推导时先在这里犹豫了一下（担心 $\E[RU]$ 需要额外信息），
建议补这半句，可以省掉读者一次回溯。

\textbf{S4（引理 8.2，实质性建议）。} 常数 $3/2$ 不紧：真实上确界为
$\sup_{[0,1]}z_*^2/\sigma=1.200224$、$\sup_{[-1,0]}z_*^2/\sigma=0.615384$（\num）。
建议二选一：(a) 把 $3/2$ 改紧为例如 $6/5$（正半轴）与 $4/5$（负半轴），
这会直接放大 \peq{contact-paid} 中 $\sigma-\theta z_*^2$ 的可用份额，是\emph{最便宜}的常数改进；
(b) 若为了有理数算术的整洁而保留 $3/2$，请在文中说明这是有意的取舍，
否则读者会以为 $3/2$ 是必要的。我倾向 (a)，理由见 \S\ref{sec:slack} 的量化估计。

\textbf{S5（可复现性，建议加一小节）。} 稿件出现 $41503=7^3\cdot11^2$、$106106902139519$、
$725029049$、$42847433/5229378000000$、$65082103/20917512000$ 等大量整数系数，
但只说明了它们的\emph{性质}，没有说明它们的\emph{来源}。
从我的复核看，$f(w)$ 由"在 $-3/8$ 与 $1/2$ 处与 $|w|+w$ 双重接触、在 $1$ 处接触"的五条件
Hermite 插值唯一确定（稿件已写明这一点，很好），
而 \peq{Q8} 与 \peq{squares} 的系数看起来是在给定 $\ell=937/1400$、$\eta=1/10$、$\theta=6/25$
一类参数后配方得到的。建议加一段"系数的生成方式与机器复现"，
说明哪些是插值唯一确定的、哪些是人工/数值选取后再精确验证的。
这不是正确性问题（我已验证它们满足所述性质），但审稿人无法区分"算出来的"与"调出来的"，
而这恰恰是这类证书型论文最容易被质疑的地方。

\textbf{S6（定理 1.2 的陈述结构）。} 目前"至多 24 个支撑点"与"原 graphon 几乎处处等于二值阶梯核"
挤在同一段，后面又接了六个 displayed 等式与团条件。建议拆成 (a)(b)(c)(d) 四条编号断言，
因为后文（以及学位论文）会反复单独引用其中 (b) 与 (d)。

\textbf{S7（\S9 的两个半径阈值）。} 引理 6.3 要求 $r_+/r_-\le3$，命题 8.5 与引理 9.1 给出
$r_+/r_-\le\sqrt8$。$\sqrt8=2.8284<3$，所以相容且有富余，但稿件没有点明这一点。
建议加半句，因为读者会自然怀疑"两个阈值是否冲突"。

\subsection{残余风险}

\textbf{R1（文献，低概率但未排除）。} Ebrahimi Boroojeni 2008 年 SFU 硕士论文
\emph{Sum of the Largest Eigenvalues of Symmetric Matrices and Graphs} 的定理 3.4.1 我无法获取
（作者自查文件中记录了同一事项）。该论文处理的是 $\Lambda_k$ 的一般上界，
若其中已含 $k=3$ 的显式系数，则本稿的新颖性声明需要相应调整。
我的间接判断是风险较低，理由见 \S\ref{sec:novel}：STW（2026）命题 6.1 仍在用
$(3+\sqrt5)/4$，而 STW 是该方向最系统的近期工作，其文献网络覆盖了 Mohar 2009、
KLMPT 2026 等；若 2008 年的硕士论文已给出更好的 $\Lambda_3$ 系数，STW 不太可能沿用旧值。
但这是\emph{推断}，不是核实。建议在投稿前用校园网访问 SFU 的 thesis 库做一次确认。

\section{与已发表结果的关系与新颖性}\label{sec:novel}

\subsection{STW 命题 6.1 的证明在哪里损失，以及本稿改的正是那一处}

我把 STW（arXiv:2608.02429）命题 6.1 的完整证明抄录并核对了。其论证是：
由 Ky Fan，$\Lambda_k(G)=\max_{Q\in P_k(n)}\tr(A_GQ)$；对 $Q=(q_{ij})\in P_k(n)$，
\[
 \tr(A_GQ)=\sum_{i\ne j}a_{ij}q_{ij}\le\sum_{i\ne j}\max(0,q_{ij})
         =\frac12\Bigl(\mathbf1^TQ\mathbf1+\|Q\|_1\Bigr)-k,
\]
然后\emph{分别}放缩两项：$\mathbf1^TQ\mathbf1=\|Q\mathbf1\|_2^2\le\|\mathbf1\|_2^2=n$
（用 $\|Q\|_{\mathrm{op}}\le1$），$\|Q\|_1\le\gamma(k)n$（绝对投影常数）。于是
$\Lambda_k(G)\le\frac{1+\gamma(k)}2n-k$。取 $k=3$，$\gamma(3)=\frac{3+\sqrt5}{1+\sqrt5}=\frac{1+\sqrt5}2=\phi$，
得系数 $\frac{1+\phi}2=\frac{3+\sqrt5}4=1.309016994\ldots$，与稿件的对比基准完全一致（我已核对）。

\textbf{这就是本稿改进的机制，而且我认为它抓住了正确的那一处。}
把 $Q=YY^T$（正交列）的行记为 $y_i$，令 $X$ 均匀分布在 $\sqrt n\,y_i$ 上，则
$\frac{\|Q\|_1}n=\E|X\cdot X'|=A$，$\frac{\mathbf1^TQ\mathbf1}n=|\E X|^2=B$，$\E XX^T=I_3$。
STW 的两个放缩分别对应"$B\le1$"与"$A\le\gamma(3)=\phi$"，两者\emph{各自}都是紧的：
$B=1$ 在 $Q\mathbf1=\mathbf1$ 时取到，$A=\phi$ 在三维极值投影（$\gamma(3)$ 的达到元）上取到。
但\textbf{两者不能同时取到}。本稿做的事正是把这个"不能同时"量化：
\[
 \frac{A+B}2\le c<\frac{1+\phi}2 .
\]
所以本稿不是把 $\gamma(3)$ 改小（那是不可能的，$\gamma(3)=\phi$ 是精确值），
而是把两个独立上界换成一个联合上界。这是一条\emph{结构上不同}的路线，
不是对 STW 的常数微调。稿件在 \S1 里说"This comparison concerns the stated universal
order-only bound"，措辞谨慎，我认为是恰当的。

\textbf{一个由此得到的、对学位论文很重要的推论。} 既然 $\gamma(3)=\phi$ 精确，
那么任何沿 STW 路线的改进都\emph{必须}来自 $A$ 与 $B$ 的耦合，
而耦合的极限恰是 $\mathcal M=\sup_W\Sigma_3(W)$（稿件已证
$\frac{A+B}2\le\Sigma_3(W_X)\le\mathcal M$）。换言之：
\textbf{把 $\Lambda_3$ 的系数压到最优，等价于确定 $\mathcal M$}，
而这条等价关系把"图论问题"彻底翻译成了"graphon 变分问题"。
这是本稿在方法论上最有价值的贡献，比常数本身更值得写进学位论文。

\subsection{与其他近期工作的关系}

我核对了本地可得的三篇相关工作：
\begin{itemize}[leftmargin=1.6em,itemsep=2pt]
\item \texttt{arXiv:2604.00512}（Kumar--Liu--Monterde--Pragada--Tait，\emph{Maximum spectral sum
of graphs}）：处理 $\Lambda_2$ 与 graphon 椭圆/团克隆论证。稿件在 \S1.1 明确引用了它作为
"机制上的先行者"，措辞是"Related graphon ellipse and clique-cloning arguments for the
two-eigenvalue problem appear in \dots"，这个归属是准确的，没有夸大。
\item \texttt{arXiv:2607.23081}（Huang--Wei）与 \texttt{arXiv:2603.21181}：涉及单个 $\lambda_3$
或 $\Lambda_2$，与 $\Lambda_3$ 的系数问题不冲突。
\item STW 的主定理（定理 2.1）针对的是\emph{单个}特征值 $\lambda_k$，
其摘要中"Our bound is tight for $k\in\{2,3,4,8,24\}$"指的是 $\lambda_k$ 那条界，
\emph{不是}命题 6.1 的 $\Lambda_k$ 界。这一点必须澄清，否则容易误读成"$\Lambda_3$ 的界已紧、
本稿不可能改进"。我特意回到 STW 摘要原文确认了这个指代：紧性声明紧跟在
$\lambda_k(G)\le\frac{(k-2)\sqrt k+\sqrt{k+1}+2}{2k(k-1)}n-1$ 之后，与 $\Lambda_k$ 无关。
\end{itemize}

\textbf{结论。} 在我能核对的文献范围内，\emph{没有}已发表的、系数优于 $c=1.284673\ldots$ 的
$\Lambda_3$ 通用上界。本稿的改进是真实的、机制上新的。

\subsection{STW \S6.1 提到的"possible gap"与本稿无关}

STW 在讨论一节提出 $c_{r+1}\le\gamma(r)/(2r)$ 可能是严格不等式（其 Table 1 的数值实验提示如此）。
这里的 $c_k=\sup\lambda_k(G)/|V(G)|$ 是\emph{单个}特征值的极限常数，与本稿的 $\Lambda_3$ 是不同对象。
本稿既没有解决也没有依赖这个问题。稿件没有把它误当作自己的成果，这一点处理得当。

\subsection{新颖性声明的边界}

稿件在 \S1.1 写了两条自我限定，我认为都必要且准确，建议在答辩时主动重申：
(i) "Our separating-observation lemma is stated and proved with its specific hypotheses in
Section 4; we do not claim a new general bound for such dependencies."（即不声称改进了
Li--Liu 的一般结果）；(ii) "The contribution here is the particular joint control of these errors
for the mixed signed and absolute energy, rather than the use of polynomial majorants by itself."
（即不声称发明了多项式上界法）。这两条限定显著降低了被审稿人以"增量性"为由打回的风险，
但也意味着\textbf{本稿的新颖性完全押在"联合控制"这一件事上}——而这恰好是与学位论文主线
（求精确上界）最相关的部分，所以定位是自洽的。

\section{常数余量的量化评估}\label{sec:slack}

这一节是我认为对学位论文最有直接价值的部分。稿件诚实地声明了 $c$ 不是最优值
（"nor assert that $c$ is optimal"，以及 \S10 "Determining the optimum remains a separate
problem"），但没有量化\emph{余量藏在哪里}。我做了这个量化。

\subsection{缺口恰好等于冠军律上的证书松弛}

设 $D_*=(3+\sqrt5)/2-2\mathcal M$，$d=(3+\sqrt5)/2-2c$。稿件的下界构造（引理 2.3 的加权轮图）
给出 $\mathcal M\ge\rho_0$，对应
\[
 D_{\rm champ}=(3+\sqrt5)/2-2\rho_0=0.054897595005559\ldots ,\qquad
 d=0.048687477302058\ldots ,
\]
\[
 \boxed{\ D_{\rm champ}-d=0.006210117703501=2\,(c-\rho_0)\ }
\]
我用两种独立方式算这个差（一次由 $c,\rho_0$ 的闭式，一次由冠军律上证书各项之和），
两者在 $10^{-12}$ 内一致。所以\textbf{系数缺口 $c-\rho_0=3.105\times10^{-3}$ 完全等价于
"冠军律在 \peq{certificate-final} 中的松弛量"}，不需要任何额外解释。

\subsection{松弛的分布：八项分解}

在冠军律（$m=6$，$a=0.197271565$，$b=0.160545687$）上，\peq{certificate-final} 右端八项的数值：
\begin{center}
\small
\begin{tabular}{lrrr}
\toprule
项 & 数值 & 占总松弛 & 备注\\
\midrule
$d$（常数项） & $0.048687477302$ & — & 这就是 $c$ 的全部来源\\
$r_{\rm tr}$ & $0.003816495539$ & $\mathbf{61.46\%}$ & \textbf{最大的未开发资源}\\
$\widehat Q=Q-V/9000000$ & $0.001648008885$ & $\mathbf{26.54\%}$ & 第二大\\
$\ell\,r_{\rm reg}$ & $0.000461837632$ & $7.44\%$ & 回归余量\\
$\mathcal R_0$ & $0.000228461432$ & $3.68\%$ & \\
$r_{\rm off}/9$ & $0.000048576162$ & $0.78\%$ & 原子交叉项\\
$K^2/1125$ & $0.000005749656$ & $0.09\%$ & \\
$\frac\kappa9(\chi-\frac{C^2}{24})$ & $0.000000988398$ & $0.02\%$ & 依赖 $m\le24$ 的那一项\\
\midrule
合计 & $0.054897595006$ & $100\%$ & 残差 $8.3\times10^{-17}$\\
\bottomrule
\end{tabular}
\end{center}

\noindent\textbf{三个立刻可读出的结论。}

第一，\textbf{支撑定理（$m\le24$）对常数 $c$ 的数值贡献几乎为零}：
$\frac\kappa9(\chi-\frac{C^2}{24})$ 只占 $0.02\%$，与我在 \S\ref{sec:method} 指出的
$c=(3+\sqrt5)/4-d_1/2-1/6000000$（原子项只贡献 $1.67\times10^{-7}$）完全一致。
这有一个重要的方法论含义：\textbf{想把 $c$ 压下来，改进 \S2--\S4 的结构理论几乎无用}，
必须改 \S5--\S8 的联合矩证书。结构定理的价值在别处（它把 $\mathcal M$ 的求解归约为有限模板问题，
见 \S\ref{sec:thesis} 的 P2），而不在常数本身。

第二，\textbf{$r_{\rm tr}$ 是最大的一块，且它的来源非常具体}。由 \peq{reserve}，
$r_{\rm tr}=\E[RR'\{(8/9-z_{\rm rad})\sigma(w)+z_{\rm rad}(\sigma(w)-g(w))\}]$，
其中 $z_{\rm rad}=3(R-R')^2/(4RR')\in[0,8/9]$，$g(w)=(1-9w^2)/25$。在冠军律上：
\[
 \frac89P=0.000515860838,\qquad
 -\E[RR'z_{\rm rad}g]=0.003300634701,\qquad
 r_{\rm tr}=0.003816495539 .
\]
也就是说 $r_{\rm tr}$ 的 $87\%$ 来自 $-\E[RR'z_{\rm rad}g]$，而它之所以为正，是因为
$V-9J=-0.055011<0$（等价地 $9J-V=2V-9\|Z\|_F^2=+0.055011$）。
换句话说：\textbf{引理 6.3 的径向支付 $\frac3{50}(V-9J)$ 在冠军律上是负的，等于什么都没支付}，
于是这一份信息被整个丢进了"只知其非负"的 $r_{\rm tr}$。稿件自己写了
"The expression $V-9J$ need not be nonnegative"，说明作者知道这一点，
但目前的处理方式是把损失记账为 $r_{\rm tr}\ge0$ 而不追讨。这是\emph{最值得追讨}的一块。

第三，\textbf{$Q$ 里有一处 $15$ 倍的显式松弛}。$Q$ 的六个子项在冠军律上为
\begin{center}
\small
\begin{tabular}{lrrr}
\toprule
子项 & 数值 & 占 $Q$ & 占总松弛\\
\midrule
$Q_3=(\frac{183}{200}-\mathrm{BIG}\,q^2)\,e$ & $0.000948524442$ & $57.56\%$ & $15.27\%$\\
$Q_1=\frac{13}{200}(\sqrt V-\frac{937}{273}q\sqrt e)^2$ & $0.000587454565$ & $35.65\%$ & $9.46\%$\\
$Q_2=\frac{19}{50}(\sqrt J-\frac{2811}{2128}q\sqrt e)^2$ & $0.000058257751$ & $3.54\%$ & $0.94\%$\\
$Q_6=\frac{65082103}{20917512000}J$ & $0.000034183572$ & $2.07\%$ & $0.55\%$\\
$Q_4=\frac{148831}{181104000}e$ & $0.000019233982$ & $1.17\%$ & $0.31\%$\\
$Q_5=\frac{42847433}{5229378000000}V$ & $0.000000359448$ & $0.02\%$ & $0.01\%$\\
\bottomrule
\end{tabular}
\end{center}
$Q_3$ 的系数 $\frac{183}{200}-\mathrm{BIG}\,q^2$ 在冠军律上是 $0.040527021596$，
但证书只能用它在整个小分支 $q^2\le20/57$ 上的下界
$\frac{183}{200}-\frac{20}{57}\mathrm{BIG}=\frac{53754367}{19868284800}=0.002705536363$，
\textbf{相差 $14.98$ 倍}。这是全稿最集中、最容易回收的一处松弛：
它不需要新想法，只需要把 $q^2$ 保留为变量、与同样依赖 $q^2$ 的其他项联合配方，
而不是先对 $q^2$ 取最坏值。

\subsection{小分支的窗口极窄，这决定了改进的难度结构}

$D_*<1/20$ 等价于 $\mathcal M>1.284016994374948$。而
$\rho_0=1.281568196872168<1.284017$，$c=1.284673255723918>1.284017$。于是：
\begin{itemize}[leftmargin=1.6em,itemsep=2pt]
\item 若 $\mathcal M\le1.2840170$，则\emph{互补分支}已经足够，$c$ 可立刻改进到 $1.2840170$
（改进 $6.56\times10^{-4}$），完全不需要动 \S5--\S8。
\item 但稿件不能这样写，因为它无法先验地知道 $\mathcal M\le1.2840170$；两个分支必须都覆盖。
\item 小分支的窗口是 $\mathcal M\in(1.2840170,\,1.2846733]$，宽度仅 $6.56\times10^{-4}$。
\item 我的数值搜索（\S\ref{sec:num}）在所有 $k\ge6$ 的原子律上都得到
$\sup\frac{A+B}2=\rho_0$，从未超过 $1.2840170$。这\emph{强烈提示}小分支实际上是空的。
\end{itemize}
\noindent 这给出一个清晰的分层：

\textbf{(i) 把 $c$ 改进到 $1.2840170$}：需要小分支给出 $D_*\ge1/20$，即从 $0.00621$ 的松弛里回收
$1/20-d=0.001312523$（占松弛的 $21.1\%$）。仅 $Q_3$ 一项的 $15$ 倍松弛就提供了 $0.000949$（$15.3\%$），
加上 $Q_1$ 的部分即可完成。\textbf{我判断这一档是可实现的}，属于"重新精确配方"而非"新想法"。

\textbf{(ii) 把 $c$ 改进到任意接近 $\rho_0$}：需要证书在冠军律上\emph{逐项取零}，
即构造一个以加权轮图为等号情形的对偶证书。这不是调常数能做到的，
它要求整个 \S5--\S8 的框架被替换成"以冠军律为极值点"的构造。
\textbf{这一档等价于证明 $\mathcal M=\rho_0$}，是本方向的真正难点。

\textbf{(iii) 证明 $\mathcal M=\rho_0$}：见 \S\ref{sec:thesis} 的 P2。

\subsection{其余可量化的松弛}

\begin{itemize}[leftmargin=1.6em,itemsep=2pt]
\item $z_*^2\le\frac32\sigma$：真实上确界 $1.200224$（$[0,1]$）与 $0.615384$（$[-1,0]$），
松弛约 $25\%$ 与 $24\%$。但它进入证书的路径是 $r_{\rm ct}\to\frac{r_{\rm ct}}9$，
在冠军律上只占松弛的 $0.80\%$（$r_{\rm off}/9$）$+0.02\%$（原子项）。
所以\textbf{收紧 $\frac32$ 的直接收益很小}，除非同时把 \peq{Cstar8} 中的除数 $125$ 一起重新优化
（$125$ 与 $\frac32$ 是耦合的：$\theta$ 变小允许用更小的除数，付出更多 $\sigma$）。
这一项我在 \S\ref{sec:issues} 的 S4 里建议改紧，理由主要是\emph{表述诚实性}而非收益。
\item 半径比阈值：引理 6.3 用 $r_+/r_-\le3$，从而 $(R-R')^2\le\frac43RR'$、
$P\ge\frac34\E[(R-R')^2\sigma]$。引理 9.1 实际能提供的是 $\sqrt8=2.8284<3$，
对应 $\frac{(R-R')^2}{RR'}\le\sqrt8+\frac1{\sqrt8}-2=\frac{9\sqrt2}4-2\le\frac{32}{27}$
（稿件 \S8 已经在用这个更紧的 $\frac{32}{27}$，它正是 $z_{\rm rad}\le\frac34\cdot\frac{32}{27}=\frac89$
的来源）。若把引理 6.3 的支付系数也从 $\frac34$ 提到 $\frac{27}{32}$，增益为 $12.5\%$ 的 $P$，
而 $P$ 在冠军律上只有 $0.000580$，即最多回收 $7\times10^{-5}$，占总松弛 $1.2\%$。收益很小。
值得注意的是这里存在\emph{内部不一致的观感}：\S6 假设比值 $\le3$、\S8 假设 $\le\sqrt8$，
两个假设都对，但读者会以为 \S6 的 $\frac34$ 是被 \S8 的 $\frac{32}{27}$ 改进了而没有用上。
建议在 \S6 就写成"比值 $\le\kappa$ 时系数为 $1/\kappa'$"的一般形式，把 $\kappa=3$ 与
$\kappa=\sqrt8$ 作为两个代入值。
\item $\alpha\ge26/105$、$\beta\ge6/7$、$\rho\ge1/2$：这三个有理下界的损失分别是
$\alpha-26/105=3.91\times10^{-4}$、$\beta-6/7=1.69\times10^{-2}$、$\rho-1/2=3.38\times10^{-2}$，
但它们只用于推出 \peq{smallbranch} 的\emph{阈值}（$e<1/20$ 等），不直接进入 $d$，
所以损失是间接的。真正进入 $d$ 的有理化损失只来自
$\sqrt5<161/72$（误差 $4.31\times10^{-5}$）与 $\sqrt3>433/250$（误差 $5.08\times10^{-5}$）：
它们给出 $d_0$ 的\emph{上}界，从而使 $\frac1{20}-d_0$ 的下界比真值小
$3.89\times10^{-5}$。（另一方向的 $d_0-\frac7{144}$ 则用 $682/305<\sqrt5$、$\sqrt3<97/56$。）
这个损失比 $0.00621$ 的结构性松弛小两个数量级。
\textbf{结论：把有理近似取得更精确，几乎不能改进 $c$}；这条路不值得走。
\end{itemize}

\subsection{一句话总结}

$c$ 与 $\rho_0$ 之间的 $3.1\times10^{-3}$，有 $61\%$ 来自"径向支付为负、损失被记为 $r_{\rm tr}$"，
有 $27\%$ 来自 $Q$（其中一半是 $Q_3$ 的 $15$ 倍最坏值放缩），
剩下 $12\%$ 分散在回归余量与 $\mathcal R_0$。
改进的性价比排序是：$Q_3$ 重新配方 $\gg$ $r_{\rm tr}$ 追讨 $\gg$ 其他。

\section{对学位论文主线的影响与建议}\label{sec:thesis}

学位论文的主线是"在 STW 命题 6.1 上改进 $\Lambda_3$ 的上界，并把求一个精确上界作为主攻方向，
下界作为副线"。本节的判断建立在前两节的量化数据上，我尽量把"能做的"与"做不动的"分开说清楚。

\subsection{必须先接受的一个定位修正}

\textbf{目前手上的是一个\emph{改进的显式上界}，不是一个\emph{精确上界}。}
精确上界意味着确定 $\mathcal M=\sup_W\Sigma_3(W)$（等价地，确定
$\sup\frac12(\E|X\cdot X'|+|\E X|^2)$），而稿件给出的是区间
\[
 \rho_0=\frac{9+\sqrt5+2\sqrt{35}}{18}=1.281568196872\le\mathcal M\le c=1.284673255724 .
\]
我的数值搜索强烈提示 $\mathcal M=\rho_0$（\S\ref{sec:num}：$k=6$ 到 $20$ 的全局最大值
\emph{精确}落在 $\rho_0$，从未超过），但这是 \num，不是结果。
所以学位论文的表述必须区分三个层次：
\begin{center}
\small
\begin{tabular}{@{}p{0.20\textwidth}p{0.34\textwidth}p{0.34\textwidth}@{}}
\toprule
层次 & 内容 & 状态\\
\midrule
A. 改进的显式上界 & $\Lambda_3(G)\le cn-3$，$c<(3+\sqrt5)/4$ & \textbf{已完成}（依赖 G1、G2 修补）\\
B. 结构定理 & 每个极大 graphon 的典范律 $\le24$ 个原子 & \textbf{已完成}（同上）\\
C. 精确上界 & $\mathcal M=\rho_0$，从而 $\Lambda_3(G)\le\rho_0 n-3+o(n)$ & \textbf{未解决}，且是真正的难点\\
\bottomrule
\end{tabular}
\end{center}
\noindent 我的建议是：学位论文把 A$+$B 作为\emph{已完成的主定理}写足，
把 C 作为"主线的第二阶段目标"单列一章，并给出本报告 \S\ref{sec:slack} 的量化分析与部分进展。
这样即使 C 在答辩前没有完成，论文也是完整且有分量的；
而如果把 C 写成主线却只交出一个非最优的 $c$，答辩时会被直接问住。

\subsection{当前框架的天花板：一个必须尽早知道的结构性事实}

这是一个必须尽早知道的结构性事实。稿件的证明是\emph{二分支}的：
\begin{itemize}[leftmargin=1.6em,itemsep=1pt]
\item 大分支 $D_*\ge1/20$：直接给 $\mathcal M\le\frac{(3+\sqrt5)/2-1/20}{2}
=\frac{3+\sqrt5}4-\frac1{40}=\frac{29+10\sqrt5}{40}=1.2840169943749\ldots$。
\item 小分支 $D_*<1/20$：用 \S5--\S8 的证书给 $\mathcal M\le c=1.2846732557239\ldots$。
\end{itemize}
两者取较坏者得 $c$。因此\textbf{在不改变分支结构的前提下，$c$ 的极限就是
$\frac{29+10\sqrt5}{40}=1.2840170$}：把小分支证书加强到"给出 $D_*\ge1/20$"，
就等于在假设 $D_*<1/20$ 下推出 $D_*\ge1/20$，即\emph{证明小分支为空}。
再往下（$<1.2840170$）就必须换掉整个分支框架。

而小分支为空这件事，数值上是成立的：$D_*<1/20\iff\frac{A+B}2>1.2840170$，
而我在全局搜索中得到的 $\sup\frac{A+B}2=\rho_0=1.2815682<1.2840170$。
冠军律本身也落在大分支（$D_*=0.0548976>0.05$，\S\ref{sec:num}）。
所以：\textbf{小分支是一个在已知构造中从未被触发的逻辑分支，
而常数 $c$ 却完全由它决定。}这是本稿最别扭、也是最值得改进的地方。

\subsection{建议的优先级排序}

\subsubsection*{P0（零风险，1--2 天，先做）}
补上 G1（4 行）与 G2（1 句）两处修补，并按 S1--S7 修改表述。
理由：这两处是\emph{承重}的（G1 撑起阈值公式与 $A+B=2\mathcal M$，G2 撑起 $m\le24$，
而 $m\le24$ 撑起 \peq{certificate-final} 中 $\chi-\frac{C^2}{24}\ge0$）。
在未修补的状态下，主定理严格来说是"有一个可修补漏洞的证明"，不能拿去投稿或答辩。
同时加一个附录"常数的生成方式与机器复现"（S5），把 \texttt{\_aud1}--\texttt{\_aud15}
的判据列成表——这会显著降低审稿人对"系数是调出来的"的怀疑。

\subsubsection*{P1（主攻第一档，中等难度，我判断可做）}
\textbf{目标}：证明小分支为空，从而
\[
 \Lambda_3(G)\le\frac{29+10\sqrt5}{40}\,n-3=1.2840169943749\ldots\,n-3 ,
\]
把 STW 的改进量从 $0.0243437$ 扩大到 $0.025$ 整（$\frac{3+\sqrt5}4-\frac{29+10\sqrt5}{40}=\frac1{40}=0.025$，
是一个非常干净的数）。

\textbf{为什么值得做}：除了常数改进本身，成功后 $c$ 的表达式会从
$\frac{106106902139519}{107310156800000}+\frac{41\sqrt3}{240}-\frac{53}{9000000}$
（一个明显是"凑出来的"数）变成 $\frac{29+10\sqrt5}{40}$（一个有闭式意义的数），
论文的可读性与可信度都会大幅提升。而且两个分支给出同一个常数，本身是自洽性的强检验。

\textbf{量化目标}：需要从当前 $0.006210$ 的证书松弛中回收
$1/20-d=0.001312523$，即 $21.1\%$。按 \S\ref{sec:slack} 的分解，
候选资源是 $r_{\rm tr}$（$61.46\%$）与 $\widehat Q$（$26.54\%$）。

\textbf{方法（按优先级）}：
\begin{enumerate}[leftmargin=1.8em,itemsep=3pt]
\item \textbf{不要先对 $q^2$ 取最坏值。} 这是最便宜的一处。当前 $Q_3$ 的系数
$\frac{183}{200}-\mathrm{BIG}\,q^2$ 被替换成它在 $q^2\le20/57$ 上的下界 $0.0027055$，
而在冠军律上它是 $0.040527$（$15$ 倍差距）。正确做法是把 $q^2$ 保留为一个变量，
把 $Q$ 看成 $(\sqrt e,\sqrt V,\sqrt J,q)$ 的二次型（系数含 $q^2$），
在小分支的约束集 $\{e<1/20,\ V<21/104,\ q^2\le20/57,\ J\le V/3\}$ 上做\emph{参数化}配方，
而不是逐点取最坏值。技术上这是一维（$q^2$）的分段有理配方，
完全可以用精确有理算术完成并验证，不需要浮点 SDP 的结论作为证明的一部分。
\item \textbf{给 $r_{\rm tr}$ 一个下界。} $r_{\rm tr}=\frac89P-\E[RR'z_{\rm rad}g]$，
其中 $-\E[RR'z_{\rm rad}g]$ 是 $r_{\rm tr}$ 的 $87\%$。它来自 $g(w)=(1-9w^2)/25$ 在
$|w|>1/3$ 上为负。要把它变成常数，需要在小分支下证明
$\E[RR'z_{\rm rad}(-g)]\ge\theta>0$。可行的切入点：小分支迫使 $A+B$ 接近 $2.618$，
这要求律高度集中（$e<1/20$ 意味着 $|\E X|^2>0.95$，即几乎全部质量在一个方向上），
而高度集中与 $z_{\rm rad}$ 的小值（$R$ 接近常值）不能同时发生——
\emph{这个"不能同时"如果能量化，就是 P1 的关键引理}。我建议先专门研究
"小分支下 $V$ 与 $e$ 的联合可行域"，这是一个二维的、可以画图分析的问题，
比直接攻证书更容易看出结构。
\item \textbf{把 \peq{smallbranch} 的阈值收紧。} 目前 $e<1/20$、$V<21/104$、$C>5/3$
都是由 $\alpha\ge26/105$、$\beta\ge6/7$、$\rho\ge1/2$ 这三个粗下界推出的。
用 $\alpha,\beta,\rho$ 的精确值（它们都在 $\mathbb Q(\sqrt5)$ 中）重推，
可以把 $e$ 与 $V$ 的上界收紧若干个百分点，从而给第 1、2 条更大的操作空间。
这一步纯机械，风险为零。
\end{enumerate}

\textbf{我的信心评估}：第 1 条与第 3 条我有较高信心（属于重新配方，不依赖新想法）；
第 2 条是真正的不确定项。综合看，我认为 P1 \emph{有实质成功机会但不是稳的}，
建议在投入前先花两天做第 3 条（零风险）并画出小分支下 $(e,V,J,q^2)$ 的可行域，
用图形判断第 2 条是否有希望，再决定是否 all-in。

\subsubsection*{P2（主攻第二档，高难度，这是"精确上界"的真正含义）}
\textbf{目标}：证明 $\mathcal M=\rho_0$，即 $\Lambda_3$ 的精确渐近上界是
$\rho_0 n-3+o(n)$，从而上下界闭合。

\textbf{为什么这不是 P1 的延续}：P1 只需回收 $21\%$ 的松弛；
P2 需要\emph{全部}回收，即构造一个在冠军律上\textbf{逐项取零}的证书。
而目前八项在冠军律上的值是
$r_{\rm tr}=3.8\times10^{-3}$、$\widehat Q=1.6\times10^{-3}$、$\ell r_{\rm reg}=4.6\times10^{-4}$、
$\mathcal R_0=2.3\times10^{-4}$、$r_{\rm off}/9=4.9\times10^{-5}$、$K^2/1125=5.7\times10^{-6}$、
$\frac\kappa9(\chi-\frac{C^2}{24})=9.9\times10^{-7}$——\textbf{没有一项为零}。
所以 P2 不是"把常数调紧"，而是要换一个以加权轮图为等号情形的新构造。
这是本质性的困难，必须如实告诉导师。

\textbf{两条可能的路线}：
\begin{enumerate}[leftmargin=1.8em,itemsep=3pt]
\item \textbf{对偶/SOS 路线}。把问题写成：在所有满足 $\E XX^T=I_3$ 的律上最小化
$D=\frac{3+\sqrt5}2-\E|X\cdot X'|-|\E X|^2$，然后构造一个对偶可行解，
其在冠军律上互补松弛全部成立。冠军律的等号条件已经完整知道
（\S\ref{sec:num}：八条结构方程全部精确成立，包括符号模式 $\mathbf1_{\{v_i\cdot v_j>0\}}=W_{ij}$
与阈值公式 $\sum_jp_jB_{ij}v_j=Lv_i$），这为"猜对偶解"提供了强约束。
具体做法是把 \peq{certificate-final} 的各个非负项换成\emph{在冠军律上消失}的项：
例如 $r_{\rm tr}$ 需要换成一个在 $9J=V$ 时消失的量，但冠军律上 $9J-V=0.055\ne0$，
所以 $r_{\rm tr}$ 这一项\emph{根本不该出现}在一个紧证书里——
这说明引理 6.3 的径向支付路线在冠军律上是"付错了方向"的，
\textbf{紧证书很可能需要放弃 \peq{payment} 这条路径}。这个判断本身就是 P2 的起点。
\item \textbf{结构刚性路线}。用定理 1.2 把 $\mathcal M$ 的求解归约为有限模板问题：
求所有满足
$\sum_ip_iv_iv_i^T=I_3$、$v_i^TLv_i=\mathcal M$、$v_{i1}>0$、$B_{ij}=\mathbf1_{\{v_i\cdot v_j>0\}}$、
$\sum_jp_jB_{ij}v_j=Lv_i$、且每个正团上 $\{v_iv_i^T\}$ 线性无关的
$m\le24$ 元系统，并证明其 $\mathcal M$ 的最大值是 $\rho_0$、唯一达到元是六原子加权轮图。
这是一个纯代数问题，原则上可判定，但 $m\le24$ 时符号模式的空间太大，
直接分类不现实。\textbf{可行的第一步是把 $24$ 降到 $6$}：
若能证明每个极大元的原子数 $\le6$（甚至证明正团结构必为 $K_1\vee C_5$），
问题规模会下降几个数量级。我注意到冠军律的正团结构正是 $K_1\vee C_5$，
而 \S4 的"$|J|\le6$"已经是 $6$，所以"整个支撑就是一个正团"是一个自然的猜想——
但这个猜想与 $\le24$（四个符号类各 $\le6$）之间还有距离，需要新的论证。
\end{enumerate}

\textbf{我的信心评估}：P2 在当前信息下\emph{不宜承诺}。
我建议学位论文把它写成"目标与部分进展"，
并把已经严格完成的部分（P0$+$P1 若能完成）作为答辩主体。

\subsubsection*{P3（副线：下界，低风险，建议完成）}
\textbf{目标}：把 $\rho_0$ 的下界从"加权 graphon 构造"升级为"有限简单图的严格渐近"。

\textbf{具体内容}：目前引理 2.3 给的是 graphon 层面的 $\Sigma_3(W)\ge\rho_0$。
对有限图，需要证明存在图列 $G_n$ 使 $\Lambda_3(G_n)/n\to\rho_0$。
我用块大小按 $a:b=0.197271565:0.160545687$ 取整的 $K_1\vee C_5$ 吹起族验证了这一点
（$n=126$ 到 $2016$，$\Lambda_3/n$ 从 $1.25770$ 单调升到 $1.28008$，\num），
但\emph{严格化}这一步稿件没有做。严格化是初等的：设块大小为
$n_0=\lfloor an\rfloor$、$n_1=\lfloor bn\rfloor$（五个），则商矩阵
$M_n=\diag(\frac{n_i}n)(A_6+I)$ 依 $O(1/n)$ 收敛到 $M=\diag(a,b,\ldots,b)(A_6+I)$，
特征值 Lipschitz 依赖于矩阵元素（Weyl 不等式给 $|\lambda_i(M_n)-\lambda_i(M)|\le\|M_n-M\|_2=O(1/n)$），
故 $\Lambda_3(G_n)=\Sigma_3(M_n)+0\cdot(\text{块内均值零空间})$，
$\Lambda_3(G_n)/n\to\rho_0$ 且误差 $O(1/n)$。

\textbf{为什么这条副线值得单独写}：它同时给出两个有用结论——
(i) 常数项 $-3$ 是\emph{不可改进}的（因为 $\tr A_G=0$ 使 $\sum_i\lambda_i=0$，
而 $\Lambda_3\ge$ …，更直接地：任何 $\Lambda_3\le cn-k$ 型界中 $k$ 的来源是 $\sum_iQ_{ii}=3$，
稿件的转移证明显示 $k=3$ 是这条路线的精确值）；
(ii) 上界与下界的差距是\emph{真实的}、随 $n$ 线性增长的，不是有限 $n$ 效应
（$n=2016$ 时余量 $+6.2634=(c-\rho_0)n$）。
这两点合起来说明\textbf{系数问题不能靠改进常数项来回避}，
从而把 P1/P2 的必要性论证清楚。这在学位论文里是很好的"问题定位"段落。

\subsection{一句话的行动建议}

先做 P0（两天，必做）；然后花两天做 P1 的第 3 条并画出小分支可行域，
据此决定 P1 投入多少；P3 与 P1 并行推进（它独立、风险低、能立刻写成论文的一节）；
P2 作为长期目标，先做"把 $24$ 降到 $6$"这一个子问题，因为它对 P2 的两条路线都有用。
\emph{方法上请坚持稿件已经示范的原则}：所有数值/优化只用于\emph{发现}常数，
最终证书必须用精确有理算术与二次域算术逐条验证——
这是本稿最可靠的部分，也是我在这次审核中唯一没有任何怀疑的部分。

\section{复现材料}\label{sec:repro}

本次审核的全部脚本保存在工作目录 \texttt{C:\textbackslash Users\textbackslash hyji11\textbackslash
Documents\textbackslash ChatGPT\textbackslash 图论\textbackslash}，环境为
\texttt{C:\textbackslash Python314\textbackslash python.exe -X utf8}（\texttt{numpy/scipy/sympy}）。
所有脚本均可独立重跑，不依赖任何中间产物。

\begingroup\small
\begin{longtable}{@{}p{0.13\textwidth}p{0.52\textwidth}p{0.24\textwidth}@{}}
\toprule
脚本 & 内容 & 结果\\
\midrule
\endfirsthead
\toprule
脚本 & 内容 & 结果\\
\midrule
\endhead
\texttt{\_aud1.py} & 常数 $c$、$\rho_0$、STW 系数；$c\in\mathbb Q(\sqrt3)$；$d,d_0,\ell_3,\ell_4,\gamma_4$；
$\sigma$ 的三条因式分解；两处精确配方；Hermite 恒等式；$g_3,g_4$ 的范数 & $23$ 断言，通过 $23$\\
\texttt{\_aud2.py} & 引理 2.2--2.5、\S5 回归、\S7--\S8 配方的逐式核对 & 约 $67$ 断言，通过 $67$\\
\texttt{\_aud3.py} & 四个有理下界证书（含用不同 $\sqrt5,\sqrt3$ 有理近似的对照实验）；
$p_*$ 的配方；$c$ 的两个恒等式 & 全部精确成立\\
\texttt{\_aud4.py}, \texttt{\_aud9.py} & 冠军律上 $16$ 条证书恒等式的残差；$4000$ 个随机律压力测试 &
残差 $\le1.7\times10^{-16}$；违例 $0$\\
\texttt{\_aud5.py} & $k=4$--$9$ 的 $\frac{A+B}2$ 全局最大化（$60$ 组初值/档） &
$k\ge6$ 恰为 $\rho_0$；$k=10$ 处因 \texttt{exp} 溢出中断\\
\texttt{\_aud6.py} & 冠军律结构方程；$n=3$--$6$ 全部 $33864$ 个图穷举；$18$ 个结构/随机图 &
无违例；\textbf{含 2 处 bug，见下}\\
\texttt{\_aud7.py} & $k=10,12,14,16,20$ 的全局最大化（稳定参数化，$70$ 组初值/档） &
全部恰为 $\rho_0$\\
\texttt{\_aud8.py} & 冠军律正团的线性无关性、$|J|\le6$ & 通过\\
\texttt{\_aud10.py}, \texttt{\_aud11.py} & $2\times10^6+1$ 点网格扫描 $\sigma$、$z_*^2/\sigma$、$h$、$q_*$、$p_*$；
$\chi$ 的配方恒等式 & 全部符合；$z_*^2/\sigma$ 的真实上确界 $1.2002/0.6154$\\
\texttt{\_aud12.py} & $K_1\vee C_5$ 吹起族，$n=126,\ldots,2016$ & $\Lambda_3/n\to\rho_0$\\
\texttt{\_aud13.py} & \textbf{本次审核新增}：修正 \texttt{\_aud9.py} 的 \texttt{K2} 命名错误；
张量恒等式自检（$\operatorname{ct}\mathcal L_4(Z)=Z$ 等 $6$ 条） & \texttt{budget4} 残差降到 $2\times10^{-17}$\\
\texttt{\_aud14.py} & \textbf{本次审核新增}：冠军律上 \peq{certificate-final} 八项与 $Q$ 六个子项的精确分解；
$r_{\rm tr}$ 的来源分解；$Q_3$ 的 $15$ 倍松弛；小分支窗口 & \S\ref{sec:slack} 的全部数据\\
\texttt{\_aud15.py} & \textbf{本次审核新增}：修正 \texttt{\_aud6.py} 的两处 bug 后重验结构方程与穷举最小余量 &
行方程偏差 $1.7\times10^{-1}\to3.3\times10^{-16}$；最小余量 $+0.0224818$\\
\bottomrule
\end{longtable}
\endgroup

\subsection{我在自己的检验工具中发现并修正的错误}

这一小节我特意写出来，因为它关系到本报告结论的可信度。审核过程中我犯了三个错，
都不是稿件的问题：
\begin{enumerate}[leftmargin=1.8em,itemsep=2pt]
\item \texttt{\_aud9.py} 把变量 \texttt{K2} 赋成了 $\E[RR'(a+b(w^2-\frac13)-\rho(w^4-\frac15)-|w|)]$，
由上界函数恒等式这等于 $S$ 而不是 $K^2=\|T_4\|_F^2$，导致 \texttt{budget4} 报告了
$-5.8\times10^{-4}$ 的假残差。修正后为 $-2.1\times10^{-17}$。
\item \texttt{\_aud6.py} 检验行方程 $\sum_jp_jB_{ij}v_j=Lv_i$ 时用了
\texttt{einsum('ij,i,jl->il')}，按 $p_i$ 而非 $p_j$ 加权，报告了 $1.675\times10^{-1}$ 的假偏差。
修正为 \texttt{einsum('ij,j,jl->il')} 后偏差为 $3.3\times10^{-16}$。
同一脚本的穷举部分把 \texttt{worst} 初始化为 $0$，使全部负值（即全部合规情形）被丢弃，
报告"最大违例 $=0$"。修正后得到真实的最小余量 $+0.0224818$ 及其取到的图（$K_1\vee C_5$）。
\item \texttt{\_aud13.py} 初版的 $\mathcal L_4(Z)$ 实现漏掉了 $\delta_{jk}Z_{il}$ 项
（重复写了两遍 $\delta_{kl}Z_{ij}$），导致谐波分解 $\E[RR'w^4]=K^2+\frac67\|Z\|^2+\frac{C^2}5$
出现 $-3.0\times10^{-4}$ 的假偏差。修正后为 $1.1\times10^{-16}$，
并补了 $6$ 条张量恒等式自检以防同类错误。
\end{enumerate}
\noindent\textbf{方法论上的结论。} 三次错误都发生在"我用一个量替换另一个量"的地方，
而三次都被"独立重算"抓出来了。这印证了本报告采用的三轨制是必要的：
如果只做符号核对而不做数值残差检验，第 1、3 类错误不会被发现；
如果只做数值检验而不做手推，第 2 类错误会被误认为是稿件的错误。
\textbf{也正因如此，我对"稿件的 \texttt{budget4}、行方程、谐波分解都精确成立"这三个结论
的信心，高于对任何单次计算结果的信心。}

\subsection{未做的事}

为诚实起见，列出本次审核\emph{没有}覆盖的部分：
\begin{itemize}[leftmargin=1.6em,itemsep=2pt]
\item 没有复现稿件整数系数的\emph{搜索过程}（只验证了它们满足所述性质），见 \S\ref{sec:method} 局限 1。
\item 没有检索 2026 年 10 月之后的新预印本，也没有获取 Ebrahimi 2008（R1）。
\item 没有对 $n\ge7$ 的图做穷举（$n=7$ 有 $1048576$ 个图，可行但本次未做；
$n\ge8$ 不可穷举）。
\item 没有验证投稿包 \texttt{top3\_submission\_sources.zip} 与 \texttt{.tex} 源文件的逐字节一致性，
只核对了文件清单与主文件的编译产物。
\item 全局极值搜索不能排除存在某个未被 $60$--$70$ 组初值覆盖的、给出更大 $\frac{A+B}2$ 的律。
\end{itemize}

\section*{结语}

这份稿件的主定理我认为\emph{是成立的}。它的证明链很长（$1346$ 行、$34$ 个编号结果、
两条互相咬合的分支），但每一环在我独立重推后都成立；
两处跳步（G1、G2）都能用五行以内的文字补上，且补上之后没有连带影响。
稿件在学术诚实性上做得比多数同类工作好：它明确不声称 $c$ 最优，
明确不声称改进了 Li--Liu 的一般结果，明确说明 \peq{budget4} 不作为独立资源重复计入，
明确说明 $V-9J$ 可能为负。这些自我限定我在复核中逐条验证过，\emph{都名副其实}。

我唯一的实质性保留意见是常数本身：$c$ 与稿件自己给出的下界 $\rho_0$ 相差 $3.1\times10^{-3}$，
而这个缺口的 $61\%$ 集中在一项（$r_{\rm tr}$）里、$15\%$ 集中在一处 $15$ 倍的最坏值放缩里（$Q_3$）。
也就是说，\textbf{缺口不是均匀弥散的，而是集中在两三个可指认的位置}。
这对学位论文是好消息：它意味着改进有明确的靶子，而不是要推倒重来。
但同时也要清楚：现有框架的天花板是 $\frac{29+10\sqrt5}{40}=1.2840170$，
要拿到真正的精确上界 $\rho_0$，需要一个在冠军律上逐项取零的新证书，那是另一个量级的工作。

\begin{thebibliography}{9}
\bibitem{STW}
V. Sivashankar, Q. Tang, T. Wakhare,
\emph{Graph eigenvalues and projection constants},
arXiv:2608.02429v1 [math.CO], 2026.
（命题 6.1：$\Lambda_k(G)\le\frac{1+\gamma(k)}2n-k$；$k=3$ 时 $\gamma(3)=\frac{1+\sqrt5}2$。）

\bibitem{KLMPT}
H. Kumar, L. Liu, H. Monterde, S. Pragada, M. Tait,
\emph{Maximum spectral sum of graphs},
arXiv:2604.00512v2, 2026.

\bibitem{KMMP}
H. Kumar, B. Mohar, S. A. Mojallal, S. Pragada,
\emph{Stability of maximal relative projection constants},
arXiv:2609.03200v1, 2026.

\bibitem{LiLiu}
H. Li, X. Liu,
\emph{Sharp bounds for minimal dependencies of linear-form powers},
arXiv:2606.24349v1, 2026.

\bibitem{Mohar}
B. Mohar,
\emph{On the sum of $k$ largest eigenvalues of graphs and symmetric matrices},
J. Combin. Theory Ser. B \textbf{99} (2009), 306--313.

\bibitem{DL}
B. Der\k{e}gowska, B. Lewandowska,
\emph{A simple proof of the Gr\"unbaum conjecture},
J. Funct. Anal. \textbf{285} (2023), 109950.

\bibitem{Glazyrin}
A. Glazyrin,
\emph{Moments of isotropic measures and optimal projective codes},
arXiv:1904.11159v2, 2020.

\bibitem{LS}
L. Lov\'asz, B. Szegedy,
\emph{Szemer\'edi's lemma for the analyst},
Geom. Funct. Anal. \textbf{17} (2007), 252--270.

\bibitem{Ebrahimi}
J. Ebrahimi Boroojeni,
\emph{Sum of the Largest Eigenvalues of Symmetric Matrices and Graphs},
M.Sc. thesis, Simon Fraser University, 2008.（未能获取；见 \S\ref{sec:issues} 的 R1。）

\bibitem{audited}
受审稿件：\emph{A joint moment bound for the three largest graph eigenvalues},
\texttt{top3\_joint\_moment\_bound.tex}（1346 行）/ \texttt{.pdf}（19 页），2026 年 11 月 2 日。
\end{thebibliography}

\end{document}
